% Options for packages loaded elsewhere
\PassOptionsToPackage{unicode}{hyperref}
\PassOptionsToPackage{hyphens}{url}
%
\documentclass[
]{book}
\usepackage{amsmath,amssymb}
\usepackage{iftex}
\ifPDFTeX
  \usepackage[T1]{fontenc}
  \usepackage[utf8]{inputenc}
  \usepackage{textcomp} % provide euro and other symbols
\else % if luatex or xetex
  \usepackage{unicode-math} % this also loads fontspec
  \defaultfontfeatures{Scale=MatchLowercase}
  \defaultfontfeatures[\rmfamily]{Ligatures=TeX,Scale=1}
\fi
\usepackage{lmodern}
\ifPDFTeX\else
  % xetex/luatex font selection
\fi
% Use upquote if available, for straight quotes in verbatim environments
\IfFileExists{upquote.sty}{\usepackage{upquote}}{}
\IfFileExists{microtype.sty}{% use microtype if available
  \usepackage[]{microtype}
  \UseMicrotypeSet[protrusion]{basicmath} % disable protrusion for tt fonts
}{}
\makeatletter
\@ifundefined{KOMAClassName}{% if non-KOMA class
  \IfFileExists{parskip.sty}{%
    \usepackage{parskip}
  }{% else
    \setlength{\parindent}{0pt}
    \setlength{\parskip}{6pt plus 2pt minus 1pt}}
}{% if KOMA class
  \KOMAoptions{parskip=half}}
\makeatother
\usepackage{xcolor}
\usepackage{longtable,booktabs,array}
\usepackage{calc} % for calculating minipage widths
% Correct order of tables after \paragraph or \subparagraph
\usepackage{etoolbox}
\makeatletter
\patchcmd\longtable{\par}{\if@noskipsec\mbox{}\fi\par}{}{}
\makeatother
% Allow footnotes in longtable head/foot
\IfFileExists{footnotehyper.sty}{\usepackage{footnotehyper}}{\usepackage{footnote}}
\makesavenoteenv{longtable}
\usepackage{graphicx}
\makeatletter
\def\maxwidth{\ifdim\Gin@nat@width>\linewidth\linewidth\else\Gin@nat@width\fi}
\def\maxheight{\ifdim\Gin@nat@height>\textheight\textheight\else\Gin@nat@height\fi}
\makeatother
% Scale images if necessary, so that they will not overflow the page
% margins by default, and it is still possible to overwrite the defaults
% using explicit options in \includegraphics[width, height, ...]{}
\setkeys{Gin}{width=\maxwidth,height=\maxheight,keepaspectratio}
% Set default figure placement to htbp
\makeatletter
\def\fps@figure{htbp}
\makeatother
\setlength{\emergencystretch}{3em} % prevent overfull lines
\providecommand{\tightlist}{%
  \setlength{\itemsep}{0pt}\setlength{\parskip}{0pt}}
\setcounter{secnumdepth}{5}
\usepackage{booktabs}
\usepackage{amsthm}
\makeatletter
\def\thm@space@setup{%
  \thm@preskip=8pt plus 2pt minus 4pt
  \thm@postskip=\thm@preskip
}
\makeatother

\usepackage{tcolorbox}


\newtcolorbox{blackbox}{
  colback=black,
  coltext=white,
  colframe=black,
  boxsep=5pt,
  arc=4pt}
\newtcolorbox{bonus}{
  colback=blue!15,
  colframe=blue!15,
  coltext=black!80,
  boxsep=5pt,
  arc=4pt}
\newtcolorbox{reflect}{
  colback=green!5,
  colframe=green!5,
  coltext=black!80,
  boxsep=5pt,
  arc=4pt}
\newtcolorbox{assessment}{
  colback=blue!5,
  colframe=blue!5,
  coltext=black!80,
  boxsep=5pt,
  arc=4pt}
\newtcolorbox{progress}{
  colback=purple!10,
  colframe=purple!10,
  coltext=black!80,
  boxsep=5pt,
  arc=4pt}
\newtcolorbox{video}{
  colback=yellow!5,
  colframe=yellow!5,
  coltext=black!80,
  boxsep=5pt,
  arc=4pt}
\newtcolorbox{caution}{
  colback=red!5,
  colframe=red!5,
  coltext=black!80,
  boxsep=5pt,
  arc=4pt}
\newtcolorbox{feedback}{
  colback=black!5,
  colframe=black!5,
  coltext=black!80,
  boxsep=5pt,
  arc=4pt}
\ifLuaTeX
  \usepackage{selnolig}  % disable illegal ligatures
\fi
\usepackage[]{natbib}
\bibliographystyle{apalike}
\IfFileExists{bookmark.sty}{\usepackage{bookmark}}{\usepackage{hyperref}}
\IfFileExists{xurl.sty}{\usepackage{xurl}}{} % add URL line breaks if available
\urlstyle{same}
\hypersetup{
  pdftitle={Digital Filmmaking},
  pdfauthor={Ned Vankevich},
  hidelinks,
  pdfcreator={LaTeX via pandoc}}

\title{Digital Filmmaking}
\author{Ned Vankevich}
\date{Last Updated Sep 2023}

\begin{document}
\maketitle

{
\setcounter{tocdepth}{1}
\tableofcontents
}
\hypertarget{welcome}{%
\chapter*{Welcome}\label{welcome}}
\addcontentsline{toc}{chapter}{Welcome}

\begin{center}\rule{0.5\linewidth}{0.5pt}\end{center}

This is the course book for MCOM 221: Digital Filmmaking. This book is divided into 6 units of study to help you engage with the materials. The course resources and learning activities are designed not only to help prepare you for the course assessments, but also to give you opportinities to practice various skills.

Below you will find information about how to navigate this book. Please also refer the schedule in Moodle, as well as the Asseessment section in Moodle for instructions on required readings and assignments.

\hypertarget{course-notes}{%
\section*{Course Notes}\label{course-notes}}
\addcontentsline{toc}{section}{Course Notes}

You should be reading this information in the context of a Trinity Western University course offered via Moodle. If this is not the case, then this may be an unauthorized reproduction of the course. Please contact \href{mailto:elearning@twu.ca}{\nolinkurl{elearning@twu.ca}} if you have concerns.

These notes will be your guide through the learning activities and assessment strategies necessary for you to succeed in the course, so it is important for you to engage to the best of your ability and take advantage of the resources available to you through Trinity Western University.

Assessment tasks are managed in other sections of the Moodle course, so be sure to familiarize yourself with those requirements and resources.

\hypertarget{how-this-course-is-built}{%
\section*{How this Course is Built}\label{how-this-course-is-built}}
\addcontentsline{toc}{section}{How this Course is Built}

This course is primarily designed to be completed asynchronously, meaning that there are no scheduled times or places that you are required to meet, even online. You can work according to your own schedule \emph{within the six weeks you have to complete the course}. That said, this is a full university level course and there are timelines that we strongly recommend that you meet to ensure that you are succeeding in building your knowledge through the course.

It would be to your significant disadvantage to submit everything at the end of the course.

Asynchronous courses require learners to be well-organized and self-motivated, and we have included supports for you to help you develop strong learning habits that will ensure your success.

For example, there are several self-check quizzes throughout the course. These quizzes are not graded, but they can be powerful tools for you to ensure you understand key ideas and concepts. We suggest you take each quiz without the aid of your notes and textbook and multiple times until you have mastered the content. This strategy taps into three powerful learning structures that have been shown to be highly effective.

\begin{enumerate}
\def\labelenumi{\arabic{enumi}.}
\tightlist
\item
  \textbf{Effortful recall.} By intentionally trying to recall information without external aids, you are strengthening the neural pathways in your brain that lead to building new connections between ideas. One way to make recall easier is to connect key ideas to other things that you know or have experienced. For example, you might be studying World War II, and you connect the date that Canadians participated in the D-Day operation with something else meaningful to you that happened on June 6, like maybe the date you bought your first car.
\item
  \textbf{Spaced repetition.} By spreading out your attempts on the quiz (leaving a few days between attempts) you can maximize the effects of the first strategy (effortful recall) and ensure that your second or third attempts truly reflect what you know about the topic. We suggest leaving 1-3 days between attempt 1 and 2, then 4-5 days between attempt 2 and 3. You can use a tool like Trello, Notion, or Asana (free versions), or even a task list on your phone to set up a spaced repetition schedule.
\item
  \textbf{Interleaving.} This is the practice of studying a particular topic for a relatively short period of time (maybe 30-40 mins), then switching to a different topic for the same period, before going back to the original topic. We will help build this into your learning by including items from unit 1 in your unit 2-6 quizzes. You can also practice this by taking regular breaks in your work, or even by retaking a unit 1 quiz while you are working in unit 2.
\end{enumerate}

These three strategies are very effective at helping people \emph{remember} key facts about a particular topic, an important first step in learning at the university level. However, you will be asked to do much more than just remember facts. Your ultimate goal is to develop \textbf{evaluative judgement}, or the ability for you to judge for yourself the quality of your (or your peers') responses to prompts.

The discussion forums are a key way for you to do this. We have set up the forums in such a way that you will need to present a response to any given prompt before you see other learners' responses. We strongly encourage you to use this structure to formulate your own ideas before you present them in the forum, and then to use the responses of your peers to help you evaluate your own response.

Using these self-check activities in this way is designed to help you to succeed on the course assignments, upon which your final grade will be determined. These assignments will require you to \textbf{use} the facts of the course to generate unique responses to the prompts, based on your past experiences, knowledge, and ability to evaluate the quality of your own work.

\hypertarget{how-to-navigate-this-book}{%
\subsection*{How To Navigate This Book}\label{how-to-navigate-this-book}}
\addcontentsline{toc}{subsection}{How To Navigate This Book}

To move quickly to different portions of the book, click on the appropriate chapter or section in the table of contents on the left. The buttons at the top of the page allow you to show/hide the table of contents, search the book, change font settings, download a pdf or ebook copy of this book, or get hints on various sections of the book.
\includegraphics{assets/course-intro/menu.png}

The faint left and right arrows at the sides of each page (or bottom of the page if it's narrow enough) allow you to step to the next/previous section. Here's what they look like:
\includegraphics{assets/course-intro/left_arrow.png} \includegraphics{assets/course-intro/right_arrow.png}

You can also download an offline copy of this book in various formats, such as pdf or an ebook. If you are having any accessibility or navigation issues with this book, please reach out to your instructor or our online team at \href{mailto:elearning@twu.ca}{\nolinkurl{elearning@twu.ca}}.

\hypertarget{course-units}{%
\subsection*{Course Units}\label{course-units}}
\addcontentsline{toc}{subsection}{Course Units}

This course is organized into 6 units. Each unit of the course will provide you with the following information:

\begin{itemize}
\tightlist
\item
  A general overview of the key concepts that will be addressed during the unit.\\
\item
  Specific learning outcomes and topics for the unit.\\
\item
  Learning activities to help you engage with the concepts. These often include key readings, videos, and reflective prompts.\\
\item
  The Assessment section provides details on assignments you will need to complete throughout the course to demonstrate your understanding of the course learning outcomes.
\end{itemize}

\begin{caution}
Note that assessments, including assignments and discussion posts will be submitted in Moodle. See the Assessment tab in Moodle for the assignment dropboxes.
\end{caution}

\hypertarget{course-activities}{%
\subsection*{Course Activities}\label{course-activities}}
\addcontentsline{toc}{subsection}{Course Activities}

Below is some key information on features you will see throughout the course.~

\begin{reflect}
\textbf{\emph{Learning Activity}}\\
This box will prompt you to engage in course concepts, often by viewing resources and reflecting on your experience and/or learning. Most learning activities are ungraded and are designed to help prepare you for the assessment in this course.
\end{reflect}

\begin{assessment}
\textbf{\emph{Assessment}}\\
This box will signify an assignment or discussion post you will submit in Moodle. Note that these demonstrate your understanding of the course learning outcomes. Be sure to review the grading rubrics for each assignment.
\end{assessment}

\begin{progress}
\textbf{\emph{Checking Your Learning}}\\
This box is for checking your understanding, to make sure you are ready for what follows. Ways to check your learning might include self-check quizzes or questions for discussion. These activities are not graded but are critical for you to be able to begin to develop evaluative judgement in this domain of knowledge.
\end{progress}

\begin{caution}
\textbf{\emph{Note}}\\
This box signifies key notes. It may also warn you of possible problems or pitfalls you may encounter!
\end{caution}

\hypertarget{project-excellence-and-professionalism}{%
\chapter{Project Excellence and Professionalism}\label{project-excellence-and-professionalism}}

\hypertarget{overview}{%
\section*{Overview}\label{overview}}
\addcontentsline{toc}{section}{Overview}

Welcome to MCOM 221! You are about to embark on a creative adventure. The specific goal of this journey is to create attention-grabbing and alluring short films. However, the general goal is far more valuable as you learn things about yourself you never knew. This course will equip you with the knowledge to develop transferable life skills that will help you in your career and professional life after you complete your education.

In addition, there is a larger element of this course that connects us to people from the past, present, and future. Namely, storytelling.

According to communications theorist Walter Fischer, human beings are more than \emph{Homo sapiens}, we are Homo \emph{narrans}, man the storyteller. Storytelling is hardwired into our brains and our cultures and history. Stories are how we make sense of our life and our world. This is why religion, philosophy, literature, and myth have been vital to human culture and understanding. Even science is a story---it tells us how the world works. But it cannot tell us why. Only story-based meaning can reveal why life is important, how we should live our lives, and why there is something rather than nothing. Think of the Bible. If one were to take out the stories, it would be a thin book.

Don't let these ideas scare you. This is a creative course not a philosophic one. The point here is that as you increase your storytelling skills---a vital part of this course---you will join the community of narrators throughout human history who have shared their wisdom, insight, and understanding. This knowledge serves as the foundation for building identities, developing traditions, and discovering the types of meaning that make life worth living across all cultures.

This course will focus on one aspect of narratives, namely visual storytelling. How do we use pictures (and sounds) to create meaning and emotionally connect people to ourselves and each other? This will be the deeper theme operating beneath the surface of this course. (In film terms it will be the subtext beneath the text.)

In practical terms, the better you are at communicating stories (relative to each profession), the better you will be at your job. Fasten your seat belt, you are about to find out why we have focused on the above to start this course.

\hypertarget{topics}{%
\subsection*{Topics}\label{topics}}
\addcontentsline{toc}{subsection}{Topics}

This unit is divided into the following topics:

\begin{enumerate}
\def\labelenumi{\arabic{enumi}.}
\tightlist
\item
  Course Introduction\\
\item
  Excellence and Professional and Personal Development\\
\item
  15 Tips for Doing Your Best\\
\item
  Transferable Skills
\end{enumerate}

\hypertarget{learning-outcomes}{%
\subsection*{Learning Outcomes}\label{learning-outcomes}}
\addcontentsline{toc}{subsection}{Learning Outcomes}

When you have completed this unit, you should be able to:

\begin{itemize}
\tightlist
\item
  Describe excellence and why it is important\\
\item
  Define what constitutes an excellent film project\\
\item
  Articulate the big picture of why stories and creativity are important\\
\item
  Determine potential transferable life skills\\
\item
  Self-assess your strengths and weaknesses\\
\item
  Plan what you want to focus on during the course in terms of your professional and personal growth.
\end{itemize}

\hypertarget{activity-checklist}{%
\subsection*{Activity Checklist}\label{activity-checklist}}
\addcontentsline{toc}{subsection}{Activity Checklist}

\begin{reflect}
Here is a checklist of learning activities you will benefit from in completing this unit. You may find it useful for planning your work.

\begin{itemize}
\tightlist
\item
  The Importance of the Film Journal: Write your first two Film Journal entries for this unit.
\item
  Read and Reflect: Read the Introduction which sets up the course ahead and Chapter One ``Project Excellence and Professionalism'' in the course textbook.
\item
  Reflect on the 15 Tips for Doing Your Best
\end{itemize}

\emph{Note that these learning activities are ungraded, but they are designed to help you succeed in your assessments in this course.}
\end{reflect}

\hypertarget{assessment}{%
\subsection*{Assessment}\label{assessment}}
\addcontentsline{toc}{subsection}{Assessment}

See the Assessment section in Moodle for assignment details and due dates.

\hypertarget{resources}{%
\subsection*{Resources}\label{resources}}
\addcontentsline{toc}{subsection}{Resources}

Here are the resources you will need to complete this unit.

\begin{itemize}
\tightlist
\item
  Introduction and Chapter One of the course text: \emph{Digital Filmmaking: A Beginner's Guide to Mastering the Craft}, by Ned Vankevich (e-text)
\item
  Other resources will be provided in the unit.
\end{itemize}

\hypertarget{course-introduction}{%
\section{Course Introduction}\label{course-introduction}}

We begin Unit 1 marvelling at the magic of film and cinema. (Read the course text Preface.) The ability of visually-centered storytelling to cross cultures and to captivate, educate, and entertain audiences is a universal contemporary phenomenon. For instance, how a series of flickering lights and images projected at 24, 25, and 30 frames a second can create emotions and experiences we all share is a wondrous mystery.

Less mysterious is how this is done. There are techniques, rules, guidelines, and practices that can help us reach people in the ways that films, TV shows, streaming Internet programs, news, and the host of other visually-centered can make us laugh, cry, and emotionally move and engage us. The worst thing that can be said by an audience is that ``I want my five, ten, sixty, or ninety minutes back (depending on the length of what we have viewed). Learning the methods to avoid this and to engage our audience will be central to this course.

\textbf{WHAT LIES AHEAD}

This course will be divided up into ten segments or units. Each segment (or chapter) is self-contained but each section is inter-connected and vital for the others. A word to the wise---do note not skip a section, even if you think that you already know it. Each section of the course will build upon the previous one.

See the overview of the course and the ten units in the Introduction to the course text: \emph{Introduction to Digital Filmmaking: A Beginner's Guide to Mastering the Craft}.

Of special note, Unit Ten and Chapter Ten will be a summary and a celebration of what you have learned during the course. It will also be a time for awards to be handed out for outstanding and excellent work such as Awards for Best film, Best Director, Best Cinematography, Best Editing., Best Story, and Most Imaginative and Most Improved Filmmaker, etc. This is noted here to give you a goal to work towards: namely, to win one of the top awards.

As can be gleaned from the above, this course will start with the most basic elements and proceed to creating a short film with a strong beginning, middle, and end. As you move through the course keep the mindset that this will be a fun, enjoyable, and exciting adventure.

\hypertarget{activity-the-importance-of-the-film-journal}{%
\subsection*{Activity: The Importance of the Film Journal}\label{activity-the-importance-of-the-film-journal}}
\addcontentsline{toc}{subsection}{Activity: The Importance of the Film Journal}

\begin{reflect}
During this course you are encouraged to keep a ``Film Journal.'' This is important for many reasons, including helping you to brainstorm ideas, as well as ponder and process what you are learning. The journal will also help you to keep a record of your course activities, and note feedback of what does and does not work in film projects of your fellow course members. In addition, some of the exercises for this course will not be graded (for reasons that will be explained ahead) but this does not mean you will not be accountable for doing them. Your journal will be submitted at the end of the course and will play a role in the grade you receive. In short, students who engage with the course well will in general receive higher marks.

Also, note that you may be asked to use your journal entries to participate in discussions, presentations, and other learning activities in the FAR Centre Facilitated Learning Labs. Please check with your facilitator about specific due dates for activities, including journal responses.

{Unit 1 Film Journal Entry:}

Based on the importance of the course journal make your first two entries for this unit:

\begin{itemize}
\tightlist
\item
  Entry One: After reading the course text Introduction, describe your best take-aways of what you have learned.
\item
  Entry Two:Study the recommended how-to-write-a-journal resources and log what you learned from them and how it will help you during the course.

  \begin{itemize}
  \tightlist
  \item
    \href{https://karenbanes.com/how-to-start-a-creative-journal/}{How to Start a Creative Journal}
  \item
    \href{https://www.youtube.com/watch?v=hUTWo7_W0lc}{How to Journal Every Day for Increased Productivity, Clarity, and Mental Health}
  \end{itemize}
\end{itemize}

\emph{Note:your journal entries do not have to be long but they should be detailed and specific enough to indicate that you have engaged the course material and projects.}
\end{reflect}

\hypertarget{excellence-and-professional-and-personal-development}{%
\section{Excellence and Professional and Personal Development}\label{excellence-and-professional-and-personal-development}}

This course will focus on the classical approach to filmmaking, that is, making motivated film stories and executing them in ways that help audiences to fully engage the story and characters without unnecessary distractions. In doing so, the course has a lofty goal: to aim for excellence. Excellence means a high standard of being good. This is of course relative to being an Introductory course with students who are just beginning their filmmaking journey, or are taking this course as an elective. But like any course, if you shoot for an A you will do better than settling for average and mediocrity.

Excellence is important in filmmaking because our competition is great. Film is a public medium. It is designed to be shown to audiences and the quality of our projects, like many things we do in life, will be judged whether we want them to be or not.

One of the transferable life skills of this course is to help us be more professional. Thus the more we shoot for excellence, the more we will develop abilities that will serve us better in our careers ahead.

At the same time we must not confuse excellence with perfection. Our beginning films and projects will have lots of mistakes and things that do not work. This should not deter us from our goal of excellence. All great artists, athletes, business leaders, and a host of other professional make mistakes. The key is to learn from them and to keep improving our craft.

HELPFUL HINT: For insight into how great artists, athletes, and performers use mistakes to grow in their art and craft see the following book by Daniel Coyle: \href{http://danielcoyle.com/the-talent-code/}{THE TALENT CODE: GREATNESS ISN'T BORN. IT'S GROWN. HERE'S HOW}.

Our goal then is to learn how to seek to produce visually-driven, digital film stories that are technically and artistically proficient, meaningfully engaging, motivated, and fraught with excellence.

As you can see excellence and professionalism go hand in hand. In order to grow and excel in a craft we need to grow and excel at being professional. Implied here is that we also need to mature and grow personally as we develop the character traits necessary for being a mature, moral, and responsible person.

\hypertarget{activity-read-and-reflect}{%
\subsection*{Activity: Read and Reflect}\label{activity-read-and-reflect}}
\addcontentsline{toc}{subsection}{Activity: Read and Reflect}

\begin{reflect}
Read Chapter One ``Project Excellence and Professionalism'' in the course textbook.

As you study and ponder what excellence and what professional and personal development mean, log in your journal why they are important and how you hope to focus on them during the course. Which character traits are your strongest and which do you need to work on? How will you do this? To help you with these entries study this resource:

\href{https://brandyourself.com/blog/guide/how-to-grow-professionally/}{How To Grow Professionally}
\end{reflect}

\hypertarget{tips-for-doing-your-best}{%
\section{15 Tips for Doing Your Best}\label{tips-for-doing-your-best}}

The following will help you to not just do well in this Film project, but will make it fun and more enjoyable and hopefully be one of your favourite courses ever:

\begin{enumerate}
\def\labelenumi{\arabic{enumi}.}
\item
  \textbf{\emph{Keep your eye on the prize.}} Creating an excellent project and growing as a professional and better person through the process.
\item
  \textbf{\emph{Meeting Deadlines.}} Being on time is vital in any profession, especially in filmmaking where the cost of feature film shoots can be thousands of dollars per hour. Plan your time and schedule accordingly so that you make your due date. Failure to do so will result in reduced marks on your assignment.
\item
  \textbf{\emph{Avoiding Distractions.}} In our Internet-cell phone era we are inundated and surrounded by 24/7 distractions. If we succumb to them we will not meet our deadlines and thus delay our professional and personal development. Learning to discern what is vital and what is the tyranny of the unimportant and unnecessary urgent is a critical skill to develop today. Focus on it and you will go further in the course. (\textbf{Tip:} \emph{turn off all notifications on your cell phone. In fact, turn off or put your cell phone away while you are working on your Film course. You will find that this will help you to focus on and finish your assignments quicker and with more creativity}).
\item
  \textbf{\emph{Focussing on the positive}}. Most of the assignments ahead will be challenging. Cultivate a ``can do attitude.'' For context think about your taking a piano lesson course. You would not expect to sit down and immediately play chords and songs. You would have to learn to play the notes, then chords, then songs over a period of time. Filmmaking is similar. There are notes to learn (film shots), chords to learn (film scenes), before we can create melodies (film sequences) and songs (full films).
\item
  \textbf{\emph{It takes time and practice to master film}}. If you find you have time, or if you are truly interested in learning the world of film or becoming a filmmaker, I encourage you to go beyond the course exercises. For example, continue to find and shoot interesting visual compositions (they are all around you) during the course, or practice fascinating and more complicated motion shots and two and three person shots. The more you do this the faster you will grow in your craft. Think about musicians. They practice over and over everyday.
\item
  \textbf{\emph{Study the practical tips and helpful hints in these units}}. Pay attention to them. They will help fast-track your skills growth.
\item
  \textbf{\emph{Find and do what you love}}. During this course you will discover things you love to do and those that you do not. As a filmmaker, every skill set is vital for the final production -- but you will discover which roles you are better at and enjoy the most. Even if you do not go into filmmaking as a career this self-insight will be invaluable as you make life decisions ahead.
\item
  \textbf{\emph{Make a firm commitment to succeed}}. When we are double-minded or not sure of what we want, we will waiver and fall prey to distraction and self-sabotage. Resolve not to give up as you make your projects despite what ``Murphy's Law'' tosses your way. (\emph{To be covered later.})
\item
  \textbf{\emph{Strive for balance}}. The emphasis on excellence and professional and personal development must not make us ``lose our soul.'' Too many people get caught up in goals and ambition and end up failing in their marriage and relationships. Remember we are physical, mental, spiritual, emotional, and relational beings and must find a balance that lets us grow in each part equally. This is the key to personal flourishing.
\item
  \textbf{\emph{Create strong and healthy relationships}}. During this course you will discover a lot about yourself and others. Some people will be easy to work with and some won't. Some will be diligent and some won't. Choose your teammates wisely and if you have a winning rapport with them, you might want to work with them on the next project. One of the keys to success in the professional world is cultivating strong and fruitful relationships. Start the process here.
\item
  \textbf{\emph{Follow your intuition}}. If you've never experimented with following your deeper instincts or ``gut feelings'' try it during this course. Often this will take you to new creative places. It might not always work but you will learn a lot from doing so.
\item
  \textbf{\emph{Take risks}}. We don't grow unless we try new things and fail. They key is not to get discouraged but to shake off a loss and fight to win again.
\item
  \textbf{\emph{Practice, practice, practice. Learn, learn, learn}}. Your best competition is continually practicing and learning their craft. So should you. You need to cultivate both a strong work ethic and a smart ethic. Take the time to think things through before launching into something.
\item
  \textbf{\emph{Take responsibility for yourself and your future}}. Playing the blame game or the victim or continually making excuses will not get you ahead in life.
\item
  \textbf{\emph{Find inspiring quotes and pin them to your computer or wall}}. Most of us need encouragement especially when the going gets tough. Seek out inspiration that helps keep you going. The following strike some of the themes of this course:
\end{enumerate}

\begin{itemize}
\item
  \emph{``Success is not final;~failure is not fatal: It is the courage to continue that counts.''} - Winston S. Churchill
\item
  \emph{``It is better to fail in originality than to succeed in imitation.''} - Herman Melville
\item
  \emph{``You will become clever through your mistakes.''} - German Proverb in ``The Talent Code''
\end{itemize}

\hypertarget{activity-reflection}{%
\subsection*{Activity: Reflection}\label{activity-reflection}}
\addcontentsline{toc}{subsection}{Activity: Reflection}

\begin{reflect}
Which of the tips in this section appeal to you most? Which do you need to focus on during the course? Write your response in your journal and share with your facilitator and/or classmates which are most important and why. Find an inspiring quote and share it with your peers.
\end{reflect}

\hypertarget{summary}{%
\section*{Summary}\label{summary}}
\addcontentsline{toc}{section}{Summary}

In this first unit, you have had the opportunity to learn about what this course entails, why it is important, and the role that professional and personal development will play. You also had a chance to respond to tips of how to make the most of this course and what skills you can transfer from it to your life.

\hypertarget{assessment-1}{%
\section*{Assessment}\label{assessment-1}}
\addcontentsline{toc}{section}{Assessment}

\begin{assessment}
{Course Journal}

After completing this unit, including the learning activities, you are asked to make sure you are doing journal entries and when appropriate to share your responses with your facilitator and classmates when you meet.

Note that entries are expected after every unit. Your journal reflections are submitted at the end of the course as part of the Journal One: Personal Journal and self assessment.

\emph{See the Assessments section for more details on submitting your journal, as well as the grading criteria.}
\end{assessment}

\hypertarget{checking-your-learning}{%
\section*{Checking your Learning}\label{checking-your-learning}}
\addcontentsline{toc}{section}{Checking your Learning}

\begin{progress}
Before you move on to the next unit, you may want to check to make sure that you are able to:

\begin{itemize}
\tightlist
\item
  Describe excellence and why it is important.\\
\item
  Define what constitutes an excellent film project.\\
\item
  Articulate the big picture of why stories and creativity are important.\\
\item
  Determine potential transferable life skills.\\
\item
  Self-assessing your strengths and weaknesses.\\
\item
  Plan what you want to focus on during the course in terms of professional and personal growth.
\end{itemize}
\end{progress}

\hypertarget{the-filmmaking-process}{%
\chapter{The Filmmaking Process}\label{the-filmmaking-process}}

\hypertarget{overview-1}{%
\section*{Overview}\label{overview-1}}
\addcontentsline{toc}{section}{Overview}

Filmmaking is exciting. We get to create and challenge ourselves as we make projects that can wow and inspire people. It is also a challenging art and craft. Film as a medium incorporates many other arts such as acting (like theatre); sets and production design (like painting and architecture); rhythm and timing (like music); scripts (like literature); movement (like dance), plasticity of form (like sculpture); as well as its unique combination of these many other arts.

Given the many skills needed to make a great film over time, filmmakers have devised a system to make it easier for filmmakers. This process is so well honed that some filmmakers, like experimental filmmakers, can do it all: produce, direct, act, shoot, and edit their film. However, the larger the film production and the bigger the budget, the more complicated the process becomes and the more people need to do it.

Like in a story that has a beginning, middle and end, the filmmaking process is divided into three sections or segments: pre-production, production, and post-production. (There is a fourth component dealing with marketing and distribution, but this will not be covered in this course.) We start with an overall look at the filmmaking process. This will give us the big picture of what will lie ahead so that we do not get lost in the details of the many elements of filmmaking, and so that we can place what we will be doing and practicing in a larger framework.

\hypertarget{topics-1}{%
\subsection*{Topics}\label{topics-1}}
\addcontentsline{toc}{subsection}{Topics}

This unit is divided into the following topics:

\begin{enumerate}
\def\labelenumi{\arabic{enumi}.}
\tightlist
\item
  The Three Phases of Filmmaking.
\item
  Pre-Production, Production, and Post-Production.
\end{enumerate}

\hypertarget{learning-outcomes-1}{%
\subsection*{Learning Outcomes}\label{learning-outcomes-1}}
\addcontentsline{toc}{subsection}{Learning Outcomes}

When you have completed this unit, you should be able to:

\begin{itemize}
\tightlist
\item
  Describe the big picture of the film production process\\
\item
  Determine what is involved in each phase and why it is important\\
\item
  Capture an overview of the processes you will do during the course\\
\item
  Define why excellence should be the benchmark of each phase of filmmaking.
\end{itemize}

\hypertarget{activity-checklist-1}{%
\subsection*{Activity Checklist}\label{activity-checklist-1}}
\addcontentsline{toc}{subsection}{Activity Checklist}

\begin{reflect}
Here is a checklist of learning activities you will benefit from in completing this unit. You may find it useful for planning your work.

\begin{itemize}
\tightlist
\item
  Watch the video ``Shooting a Film- START to FINISH!''
\item
  Watch the video ``Pitching and Pre-Production'' to learn why pitching is important.
\item
  Read and study Chapter 2, and answer to the questions provided
\item
  Share in your journal what you learn from shooting a film from start to finish.
\end{itemize}

\emph{Note that these learning activities are ungraded, but they are designed to help you succeed in your assessments in this course.}
\end{reflect}

\hypertarget{assessment-2}{%
\subsection*{Assessment}\label{assessment-2}}
\addcontentsline{toc}{subsection}{Assessment}

See the Assessment section in Moodle for assignment details and due dates.

\hypertarget{resources-1}{%
\subsection*{Resources}\label{resources-1}}
\addcontentsline{toc}{subsection}{Resources}

Here are the resources you will need to complete this unit.

\begin{itemize}
\tightlist
\item
  Chapter Two of \emph{Digital Filmmaking: A Beginner's Guide to Mastering the Craft,} by Ned Vankevich (e-text)
\item
  Other online resources will be provided in the unit.
\end{itemize}

\hypertarget{the-three-phases-of-filmmaking}{%
\section{The Three Phases of Filmmaking}\label{the-three-phases-of-filmmaking}}

In Unit 1 we focused on project excellence and great storytelling and why they are important. This unit will help obtain these goals by breaking the filmmaking process down into three phases. Each phase depends upon the other and if we fail to understand and execute each one well, the other phases and the film itself will suffer.

The overall filmmaking process has not changed much over the past century. In short, it involves three phases: Pre-Production---preparing to shoot a project; Production---shooting or filming the project; Post-Production---editing or putting all of the production elements together in a unified final project.

Understanding these three phases is vital to making a great film. You will not have to master each phase, no one can because there are too many elements and skills involved depending on the complexity of the film, its story, and its execution.

\hypertarget{activity-shooting-a-film}{%
\subsection*{Activity: Shooting a Film}\label{activity-shooting-a-film}}
\addcontentsline{toc}{subsection}{Activity: Shooting a Film}

\begin{reflect}
To get a visual representation of the overall filmmaking process view the following resource:

Don't worry if you do not know all of the terms he is using such as blocking, dolly, tilt up, etc. These terms will be covered in the units ahead. For now, just take in the whole process and you will be amazed how much easier it will be to understand the later units as you explore the details of each phase.

For now, just watch the video and enjoy the process.

\textbf{Helpful Hint:} \emph{Don't multi-task during the viewing of resources. It will divide your attention and you will not absorb as much content. As with most things in this course---be in the moment and focus on what is at hand.}
\end{reflect}

\hypertarget{pre-production}{%
\section*{Pre-Production}\label{pre-production}}
\addcontentsline{toc}{section}{Pre-Production}

Like most things in life, if you do not have a strong foundation, what follows after will wobble. Proper pre-production is the foundation of an excellent film (and your up-coming exercises). The better you plan something, the better the result.

There are many elements to pre-production. When followed properly they enhance the 5 Cs of successful filmmaking (see the \textbf{Why Pre-Production is Important} section in Chapter 2 of the course text):

\begin{itemize}
\tightlist
\item
  Creativity\\
\item
  Calmness\\
\item
  Co-ordination\\
\item
  Coherence\\
\item
  Consistency
\end{itemize}

Proper pre-production involves many important general steps (see \textbf{What Proper Pre-Production Entails} section in Chapter 2 of the course text).

\begin{itemize}
\tightlist
\item
  Planning\\
\item
  Anticipation\\
\item
  Strategizing\\
\item
  Testing
\end{itemize}

As can be gleaned from above, pre-production helps ensure that a film project maximizes the creative process and minimizes chaos, confusion, and uncertainty---the enemies of a good film.

There are many steps in the prep-production process highlighted in Chapter 2. (See \textbf{Successful Steps of Pre-Production}.)

They involve:

\begin{itemize}
\tightlist
\item
  Creating a viable concept for a project\\
\item
  Knowing who the audience is and genre requirements (people who like comedy want to laugh, people who like horror want to be scared)\\
\item
  Creating a script to film\\
\item
  Creating storyboards and shot lists\\
\item
  Pitching your project to get funding, actors, crew members, etc.\\
\item
  Budgeting the cost of the film\\
\item
  Getting your actors and the right crew\\
\item
  Finding locations and props\\
\item
  Scheduling the shoot\\
\item
  Testing your gear to make sure you know how to use it and that it works properly.
\end{itemize}

As can be seen, there is a lot involved before we film. Luckily, we will start with short exercises which will take minimal pre-production and gradually build to the more detailed aspects later.

\textbf{Note:} The lion's share of this unit and Chapter Two in the course book is devoted to pre-production to emphasize how important this phase is. Most people might think production is the most important part of the filming process. In fact, each phase is. In Hollywood the development process (creating the screenplay) and pre-production for a film often takes far longer than production. Post-Production is also a longer process than production when a major film is involved.

\hypertarget{production}{%
\section*{Production}\label{production}}
\addcontentsline{toc}{section}{Production}

The production phase involves the shooting of the film, what is often referred to as ``principle photography''. (Sometimes pick-up shots, re-dos, and B-roll footage are done during post-production.) For many filmmakers this is the most exciting part of filmmaking when they get a chance to go on location and watch actors do their magic.

Like pre-production, the production phase must be attended to thoughtfully and diligently or you will not have what you need to edit your film and make your story work. This is why a large portion of this course and the course text are dedicated to this phase of filmmaking, which will only be covered briefly here.

In addition to filming the actors, locations, action sequences, etc. production also includes recording on-location audio for the project. Capturing, recording, and creating good sound, like cinematography, takes lots of time and practice.

Most of the production process takes place on real locations such as streets, schoolyards, restaurants, etc., and sets which are built for the filming. If you have the budget, filming on sets such as kitchens, living, rooms, police stations, etc. is easier than real locations. The reason for this is that you can control the lighting and use of space better and will not have extraneous and disruptive noises occurring in the environment.

Production can also include green screen and CGI (computer generated imagery) work, but these are normally for bigger budget productions since they are specialized skills and it's costly to do them well. They are also time consuming. (If you have After Effects and other special effects software skills please consider using them in this course.)

Performance is a crucial part of the filming process. Many filmmakers forget this and get so caught up in the shooting process that they overlook or miss bad acting. Don't fall into this trap. You want to work on having your actors deliver believable, motivated, non-melodramatic performances. More on this later.

\hypertarget{post-production}{%
\section*{Post-Production}\label{post-production}}
\addcontentsline{toc}{section}{Post-Production}

The final phase of the filmmaking process involves post-production where all elements of pre-production and production are ``cut'' or stitched together to create a finished film.

Like production, the post process is hands-on and labor intensive. If done digitally, it involves using software such as Adobe Premier Pro, Apple Final Cut, Avid, DaVinici Resolve, etc. to edit the footage and make transitions such as fades, dissolves, wipes, etc. (to be covered later). This can be easy but a lot of work goes into choosing the best and rights shots and editing them together.

There are four phases to cutting the picture of a film:

\begin{itemize}
\tightlist
\item
  Assembly Cut: where the shots are placed together in proper order without trimming them.\\
\item
  Rough Cut: where you begin to ``trim the fat'' from the beginning and end of shots to get a feel for how the story will flow.\\
\item
  Fine Cuts: which will involve multiple versions as you trim or cut shots and scenes that do not work. You can further hone the film of any ``excess'' fat that does not add to the effective presentation of the story and characters.\\
\item
  Final Cut: which is the ``locked'' picture version that the composer and sound design people need to add music and sound effects to precise moments in the story.
\end{itemize}

In addition to cutting the picture, post-production also entails:

\begin{itemize}
\item
  \textbf{Sound Design} where the right music and sound effects are chosen and added.
\item
  \textbf{Titling and Graphics} when the opening and ending film credits appear and any special written material that will appear on the screen such as: FIVE YEARS EARLIER.
\item
  \textbf{Test Screening} your film to those you trust to make sure the story works and is clear, and to test how an audience will respond.
\item
  \textbf{Color Correcting} where the flow of the color and light and dark nature of the images appears seamless and appropriate. This is where you correct over exposed and under expose shots or those with the wrong color temperature. (\emph{Look these terms up in the film glossary you are using during this course. Refer to Activity 1.2 in Unit 1})
\item
  \textbf{Audio Mixing} which involves finding the right levels and balance between the sound elements such as dialogue, music, sound effects, room tone, etc.
\item
  \textbf{Format Delivery}: What resolution will you use to export your film project? The wrong one can undermine the quality of your film.As can be gleaned, there is a lot to post-production but by carefully studying and practicing the units ahead it will make it easier.
\end{itemize}

If you are feeling overwhelmed and intimidated, you will not be alone. This is a lot to do and learn and this is why the bigger film projects can be done with team members.

\hypertarget{activity-reflect-and-share}{%
\subsection*{Activity: Reflect and Share}\label{activity-reflect-and-share}}
\addcontentsline{toc}{subsection}{Activity: Reflect and Share}

\begin{reflect}
Based on this unit and your reading of Chapter 2 in the course text, share in your journal why knowing the overall process is important and what you learned most from this unit.

Share also if you are feeling intimidated. Remember all but one of the graded assignments will be done in teams so you will not have to do it all alone.
\end{reflect}

\hypertarget{summary-1}{%
\section*{Summary}\label{summary-1}}
\addcontentsline{toc}{section}{Summary}

In this unit, you learned about the overall process of making a film as well as the steps involved in the following three phases:

\begin{itemize}
\tightlist
\item
  Pre-Production: this must be done before making a film to save time and help ensure quality, as you write the script, plan the production, and work out the logistics such as casting and finding locations and props.
\item
  Production: this involves shooting the film and where you use camera angles, camera motions, and the blocking of actors to reveal the meaning of your story and its characters.
\item
  Post-Production: this involves the editing and completing of your film where you cut the shots and scenes together, and includes adding titling, credits, special effects, music and sound design, and colour grading.
\end{itemize}

\hypertarget{assessment-3}{%
\section*{Assessment}\label{assessment-3}}
\addcontentsline{toc}{section}{Assessment}

\begin{assessment}
{Assignment 1: Course Journal}

After completing this unit, including the learning activities, you are asked to make sure you are doing journal entries and when appropriate to share your responses with your facilitator and classmates when you meet.

Note that entries are expected after every unit. Your journal reflections are submitted at the end of the course as part of the Journal One: Personal Journal and self assessment.

\emph{See the Assessments section for more details on submitting your journal, as well as the grading criteria.}
\end{assessment}

\hypertarget{checking-your-learning-1}{%
\section*{Checking your Learning}\label{checking-your-learning-1}}
\addcontentsline{toc}{section}{Checking your Learning}

\begin{progress}
Before you move on to the next unit, you may want to check to make sure that you are able to:

\begin{itemize}
\tightlist
\item
  Describe the big picture of the film production process\\
\item
  Determine what is involved in each phase and why it is important\\
\item
  Capture an overview of what lies ahead\\
\item
  Define why excellence should be the bench mark of each phase
\end{itemize}
\end{progress}

\hypertarget{visual-composition}{%
\chapter{Visual Composition}\label{visual-composition}}

\hypertarget{overview-2}{%
\section*{Overview}\label{overview-2}}
\addcontentsline{toc}{section}{Overview}

We will now make our theory more practical. So far, we have focused on project excellence and the three phases of filmmaking we will use to create it. In this unit we will concentrate on where it all begins visually -- the single image and how to craft it well.

This unit will also give you a chance to apply what you are learning in a creative way. Don't worry if you think you are not creative or artistic We will focus on simple elements and reveal how they can help you see the world and capture it in exciting ways. We will do so by focusing on composition---the way in which the elements in a picture or image are arranged.

Visual composition will play an important role in this course, so taking the time to learn it is well worth the investment.

\hypertarget{topics-2}{%
\subsection*{Topics}\label{topics-2}}
\addcontentsline{toc}{subsection}{Topics}

This unit is divided into the following topics:

\begin{enumerate}
\def\labelenumi{\arabic{enumi}.}
\tightlist
\item
  Elements and Principles of Composition
\item
  Photographic Themes
\end{enumerate}

\hypertarget{learning-outcomes-2}{%
\subsection*{Learning Outcomes}\label{learning-outcomes-2}}
\addcontentsline{toc}{subsection}{Learning Outcomes}

When you have completed this unit, you should be able to:

\begin{itemize}
\tightlist
\item
  Describe the elements used in visual compositions\\
\item
  Define composition principles\\
\item
  Study the works of noted photographers and apply composition elements and principles to photographs you create\\
\item
  Create photos that reveal your understanding of the chapter's content
\end{itemize}

\hypertarget{activity-checklist-2}{%
\subsection*{Activity Checklist}\label{activity-checklist-2}}
\addcontentsline{toc}{subsection}{Activity Checklist}

\begin{reflect}
Here is a checklist of learning activities you will benefit from in completing this unit. You may find it useful for planning your work.

\begin{itemize}
\tightlist
\item
  Read and study Chapter 3 of course e-text: \emph{Digital Filmmaking: A Beginner's Guide to Mastering the Craft.}
\item
  Read the article ``11 Thoughts: An Introduction to Photographic Composition'' by Todd Vorenkamp and reflect on the question provided.
\item
  Read the article ``20 Composition Techniques That Will Improve Your Photos'', and practice and apply techniques and principles of design.
\item
  Review the suggested websites to enrich your understanding of photographic themes and find images that stand out in intriguing, startling, fascinating, and emotionally moving ways.
\end{itemize}

\emph{Note that these learning activities are ungraded (unless specified), but they are designed to help you succeed in your assessments in this course.}
\end{reflect}

\hypertarget{assessment-4}{%
\subsection*{Assessment}\label{assessment-4}}
\addcontentsline{toc}{subsection}{Assessment}

\begin{itemize}
\tightlist
\item
  Film Journal\\
\item
  Visual Composition Exercise (ungraded, but part of journal entry)
\end{itemize}

\hypertarget{resources-2}{%
\subsection*{Resources}\label{resources-2}}
\addcontentsline{toc}{subsection}{Resources}

Here are the resources you will need to complete this unit:

\begin{itemize}
\tightlist
\item
  Chapter 3 of course text: \emph{Digital Filmmaking: A Beginner's Guide to Mastering the Craft,} by Ned Vankevich (e-text)
\item
  Other online resources will be provided in the unit.
\end{itemize}

\hypertarget{elements-and-principles-of-composition}{%
\section{Elements and Principles of Composition}\label{elements-and-principles-of-composition}}

{Elements of Composition}

At heart, visual composition means the intentional arrangement or conscious activity of constructing the ``ingredients'' of an image. As we grow in our filmmaking craft, we must continually be aware of what takes place within the frame (the border of a picture or a film shot). Focusing on a single image can help train us to do this.

\begin{caution}
\textbf{Important Tip}: Before continuing below, study this overview of the building blocks of composition in art: \href{http://www.artyfactory.com/art_appreciation/visual-elements/visual-elements.html}{The Visual Elements: The Building Blocks of Composition in Art}
\end{caution}

The following are some of the formal elements that make up the design and structure of an image. However, it is important to remember that subject matter and composition are linked and therefore do not get locked in to rigidly adhering to what you are learning. Many masters break the rules but they know them well before they do so.

\begin{itemize}
\tightlist
\item
  Line\\
\item
  Shape (geometric, organic, natural)\\
\item
  Colour (hues, intensities, symbolism)\\
\item
  Texture (tactile feel)\\
\item
  Value (shadows and shading)\\
\item
  Form (3-D, 2-D)\\
\item
  Space (positive/objects); (negative/space between objects)\\
\item
  Depth (foreground/mid-ground/background)
\end{itemize}

{Principles of Composition}

In addition to elements of design, there are also principles or general rules that help give structure to visual compositions and which can also lend meaning or a theme to an image. Elements are specific things whereas principles are more general. Here are four that the course text focuses on:

\textbf{1. Balance}\\
\textbf{2. Rule of Thirds}\\
\textbf{3. Repetition and Patterns}\\
\textbf{4. Combinations}

There are many other principles such as the use of triangles and frame within a frame in compositions.

\hypertarget{activity-read-and-study}{%
\subsection*{Activity: Read and Study}\label{activity-read-and-study}}
\addcontentsline{toc}{subsection}{Activity: Read and Study}

\begin{reflect}
Read and study Chapter 3. Be sure to study each of the images provided in this section and how the elements are used. This will increase your understanding of how they function in visual compositions.
Before exploring the rest of this chapter, study this essay to help guide you during the process we are exploring:\href{https://www.bhphotovideo.com/explora/photography/tips-and-solutions/11-thoughts-introduction-photographic-composition}{11 Thoughts: An Introduction to Photographic Composition}, by Todd Vorenkamp.

{Question to Consider}

Based on your study of this section, what are some observations you can share in your journal and with you facilitator and classmates?
\end{reflect}

\hypertarget{activity-readings-on-techniques-and-principles-of-composition}{%
\subsection*{Activity: Readings on Techniques and Principles of Composition}\label{activity-readings-on-techniques-and-principles-of-composition}}
\addcontentsline{toc}{subsection}{Activity: Readings on Techniques and Principles of Composition}

\begin{reflect}
Read the article \href{https://petapixel.com/photography-composition-techniques/}{28 Composition Techniques That Will Improve Your Photos}, and \href{https://www.getty.edu/education/teachers/building_lessons/principles_design.pdf}{Principles of Design}. Practice and apply techniques and principles of design.
\end{reflect}

\hypertarget{photographic-themes}{%
\section{Photographic Themes}\label{photographic-themes}}

The subjects of a visual image includes the objects, people, or items in the frame. E.g., a mountainous landscape, a portrait of twins, etc.

A more advanced approach to composition involves the use of creating a theme or deeper meaning to a composition, or where you take a topic such as clothes on clotheslines and photograph a series of images that reveal their beauty or what they say about the culture, people, or environment where they were shot.

As you grow in your composition skills you will discover how to add more depth, interest, and meaning to your photographs and images.

\hypertarget{activity-read-and-explore}{%
\subsection*{Activity: Read and Explore}\label{activity-read-and-explore}}
\addcontentsline{toc}{subsection}{Activity: Read and Explore}

\begin{reflect}
You do not have to create a theme for your photographs in this unit. However, studying websites like the following can help enrich your understanding and find images that stand out in intriguing, startling, fascinating, and emotionally moving ways.

\begin{itemize}
\tightlist
\item
  \href{https://www.moma.org/artists/1777}{Walker Evans MoMA Exhibit}
\item
  The outstanding Depression Era works of Dorothea Lange are also rich with theme. See \href{http://www.historyplace.com/unitedstates/lange/}{The History Place Dorothea Lange}
\item
  \href{https://www.boredpanda.com/top-10-photographers-for-travel-portraits/?utm_source=google\&utm_medium=organic\&utm_campaign=organic}{Top 10 Most Famous Portrait Photographers In The World}
\end{itemize}

\textbf{\emph{Helpful Hint:}} \emph{The photographs of Walker Evans and Dorethea Lange cited above are almost one-hundred years old. But they reveal the way their heartfelt themes and beauty are universal and timeless---the goal of excellent photos.}

Log in your journal and share with a peer in your course what you learned about photographic themes.
\end{reflect}

\hypertarget{summary-2}{%
\section*{Summary}\label{summary-2}}
\addcontentsline{toc}{section}{Summary}

In this unit, you learned about the use of elements, principles, and theme in visual composition and photographic design.

Composition and subject matter are different but intimately related. We can have a simple subject with little composition. A white ball on a black table. Or we can have a complex composition with little subject matter. The bokeh (blurred or out of focus) of rain splashes on a glass window.

We have emphasized composition here because well composed shots are a foundational building block of a good film. As you will learn, if the head space of your shots is too much or too little, or there is a lack of balance of how people are arranged, or if you do not give enough leading entry space in a jogging shot, it will detract from the effect of your film. (There are exceptions to these rules which will be addressed later.) Too many flaws will mark you as an untrained or sloppy filmmaker - something we want to avoid.

\hypertarget{assessment-5}{%
\section*{Assessment}\label{assessment-5}}
\addcontentsline{toc}{section}{Assessment}

\begin{assessment}
{Course Journal}

After completing this unit, including the learning activities, you are asked to make sure you are doing journal entries and when appropriate to share your responses with your facilitator and classmates when you meet.

Note that entries are expected after every unit. Your journal reflections are submitted at the end of the course as part of the Journal One: Personal Journal and self assessment.

\emph{See the Assessments section for more details on submitting your journal, as well as the grading criteria.}

{Composition Exercise}

\emph{This activity is ungraded, but part of journal entry}

Congratulations! You have earned the right to shoot your first assignment for the course.

After you have done some composition research and explored great photographers and photos and their compositional techniques and aesthetics, review Chapter 3 and go out and find and create five outstanding visual compositions with your camera (cell phone, DSLR, etc.). Share these with your family, friends, and peers and get their feedback about what worked and did not work. When you have found your best images upload them to the course online folder where the class assignments will be stored.

Remember, the name of the game is to learn not to be perfect.bRefer to the specific guidelines in the assessment section. Be sure to study them carefully.

\emph{See the Assessments section for more details on your composition, as well as the grading criteria.}
\end{assessment}

\hypertarget{checking-your-learning-2}{%
\section*{Checking your Learning}\label{checking-your-learning-2}}
\addcontentsline{toc}{section}{Checking your Learning}

\begin{progress}
Before you move on to the next unit, you may want to check to make sure that you are able to:

\begin{itemize}
\tightlist
\item
  Describe the elements used in visual compositions.
\item
  Define composition principles.
\item
  Study the works of noted photographers and apply composition elements and principles to photographs you create.
\item
  Create photographs that reveal your understanding of the chapter's content.
\end{itemize}
\end{progress}

\hypertarget{cinematic-motion}{%
\chapter{Cinematic Motion}\label{cinematic-motion}}

\hypertarget{overview-3}{%
\section*{Overview}\label{overview-3}}
\addcontentsline{toc}{section}{Overview}

In Unit 3 you created visual compositions with your camera. We hope you found it exciting to go out and explore your world through a lens and design a creative still image.

In this Unit you begin your journey into the world of motion pictures. At heart, film (and video) is a temporal medium. It involves a series of images that unfold over time, most often at the rate of 24, 25, or 30 frames per second. This rate of frame-flow gives film its sense of moving pictures where a series of single frame images of a galloping horse look like the horse is actually running. This is why film is considered a kinetic medium where moving people and objects take center stage.

We began with still photos in the previous unit because they are easier to control and to frame. In this unit we are going to add motion to what we compositionally frame. This will bring a lot more variables to image creation and make it more exciting for many people. As you focus on motion you must not be tempted to forget the lessons learned in the previous unit. Many of the same composition elements and principles will apply and the wise student will refresh his or her self with Unit 3 before engaging the activities for this unit.

\hypertarget{topics-3}{%
\subsection*{Topics}\label{topics-3}}
\addcontentsline{toc}{subsection}{Topics}

This unit is divided into the following topics:

\begin{enumerate}
\def\labelenumi{\arabic{enumi}.}
\tightlist
\item
  Types of Cinematic Motion\\
\item
  Motivating Camera Movement
\end{enumerate}

\hypertarget{learning-outcomes-3}{%
\subsection*{Learning Outcomes}\label{learning-outcomes-3}}
\addcontentsline{toc}{subsection}{Learning Outcomes}

When you have completed this unit, you should be able to:

\begin{itemize}
\tightlist
\item
  Describe the various types of film motion shots
\item
  Contrast the various types of film motion shots
\item
  Determine when and how to use cinematic motion
\item
  Create effective cinematic motion shots
\end{itemize}

\hypertarget{activity-checklist-3}{%
\subsection*{Activity Checklist}\label{activity-checklist-3}}
\addcontentsline{toc}{subsection}{Activity Checklist}

\begin{reflect}
Here is a checklist of learning activities you will benefit from in completing this unit. You may find it useful for planning your work.

\begin{itemize}
\tightlist
\item
  Read and Study: Read and study Chapter 4 of course e-text: \emph{Digital Filmmaking: A Beginner's Guide to Mastering the Craft.}
\item
  Motion Shots Practice: Practice doing as many of the main types of motion shots as you can. Discuss what you did and why with your peers.
\item
  Proper Motivation Practice: Find someone to film and practice each of the cinematic motion techniques with your cell phone.
\end{itemize}

\emph{Note that these learning activities are ungraded, unless specifie, but they are designed to help you succeed in your assessments in this course.}
\end{reflect}

\hypertarget{assessment-6}{%
\subsection*{Assessment}\label{assessment-6}}
\addcontentsline{toc}{subsection}{Assessment}

See the Assessment section in Moodle for assignment details and due dates.

\hypertarget{resources-3}{%
\subsection*{Resources}\label{resources-3}}
\addcontentsline{toc}{subsection}{Resources}

Here are the resources you will need to complete this unit:

\begin{itemize}
\tightlist
\item
  Chapter Four of the course text: \emph{Digital Filmmaking: A Beginner's Guide to Mastering the Craft}, by Ned Vankevich
\item
  Other online resources will be provided in the course text and unit.
\end{itemize}

\hypertarget{types-of-cinematic-motion-shots}{%
\section{Types of Cinematic Motion Shots}\label{types-of-cinematic-motion-shots}}

In Unit 3 we explored how excellent or effective still (non-moving) visual compositions are made. The elements and principles of such images apply to filmmaking too. However, film and video add new variables and techniques given that they deal with moving or motion pictures.

However, as we explore film motion, we must remember that film and video involve a series of still images being projected and that persistence of vision makes them appear connected. (See the first Helpful Hint in the course text Chapter 4.) This is why we need to incorporate what we have learned about visual composition as we create our ``moving'' pictures.

The main types of motion shots include:

\begin{itemize}
\tightlist
\item
  ZOOMS\\
\item
  PAN SHOTS\\
\item
  SWISH OR WHIP PAN\\
\item
  TILT SHOTS\\
\item
  DOLLY\\
\item
  DOLLY ZOOMS\\
\item
  TRACKING\\
\item
  ARCING\\
\item
  FOLLOW SHOTS\\
\item
  CRANE SHOTS\\
\item
  360-DEGREE TRACKING SHOTS\\
\item
  SLIDERS\\
\item
  GOPRO\\
\item
  DRONES\\
\item
  SHAKY CAMERA
\end{itemize}

\hypertarget{activity-motion-shots-reading}{%
\subsection*{Activity: Motion Shots Reading}\label{activity-motion-shots-reading}}
\addcontentsline{toc}{subsection}{Activity: Motion Shots Reading}

\begin{reflect}
Study Chapter 4 of course e-text: \emph{Digital Filmmaking: A Beginner's Guide to Mastering the Craft.} This chapter will help you understand the different types of motion shots. Be sure to take notes and begin logging and identifying shot ideas you want to try and practice. Actively taking notes as you read is a great way to absorb the material.
\end{reflect}

\hypertarget{activity-motion-shots-practice}{%
\subsection*{Activity: Motion Shots Practice}\label{activity-motion-shots-practice}}
\addcontentsline{toc}{subsection}{Activity: Motion Shots Practice}

\begin{reflect}
Study the definitions of the main types of motion shots and how to do them. Then practice doing as many of them as you can. (Obviously you cannot do a GoPro or drone shot without that specific equipment, but you might be imaginative and find an alternative to do a similar type of shot.)

Write in your journal your observations about each type of shot and what you learned from doing them.
\end{reflect}

\hypertarget{activity-watch-and-learn}{%
\subsection*{Activity: Watch and Learn}\label{activity-watch-and-learn}}
\addcontentsline{toc}{subsection}{Activity: Watch and Learn}

\begin{reflect}
Before moving on to the next topic, watch the video \href{https://www.youtube.com/watch?v=h2c3JZ6X3f8}{``5 Brilliant Moments of Camera Movement''} to augment your understanding of camera movement

\begin{center}\rule{0.5\linewidth}{0.5pt}\end{center}

\textbf{Note:} \emph{Do not try doing zoom shots with your cell camera unless you have a special app that makes it smooth. Squeezing your finger in and out to move closer or further away from your subject will result in a shaky shot. Our goal is to disguise movement, not to draw attention to it.}
\end{reflect}

\hypertarget{properly-motivated-camera-moving-shots}{%
\section{Properly Motivated Camera Moving Shots}\label{properly-motivated-camera-moving-shots}}

A classical approach to filmmaking involves using shot and techniques that are motivated by the characters and story and which do not unnecessarily distract or pull the audience out of experiencing an event in your film. Put another way, if you do something that is jarring or not done well such as a shaky zoom shot it will draw attention to itself and distract the audience and prevent them from engaging the flow of your scene. What this means will become apparent as you progress in the course.

The following sections will help you have properly motivated and framed movement within the frame.

\begin{itemize}
\tightlist
\item
  Properly Motivated Movement Shots
\item
  Proper Placement in the Frame
\item
  Motivated Unmotivated Camera Movement
\end{itemize}

\hypertarget{activity-proper-motivation-exercise-and-practice}{%
\subsection*{Activity: Proper Motivation Exercise and Practice}\label{activity-proper-motivation-exercise-and-practice}}
\addcontentsline{toc}{subsection}{Activity: Proper Motivation Exercise and Practice}

\begin{reflect}
After you have viewed the resources and studied the properly motivated camera movements, find someone to film and practice each of the techniques with your cell phone. For inspiration and how-to-do techniques watch the video \href{https://www.youtube.com/watch?v=6_p93J3OwfU}{``8 Cinematic Camera Moves For Video''}
\end{reflect}

\hypertarget{summary-3}{%
\section*{Summary}\label{summary-3}}
\addcontentsline{toc}{section}{Summary}

In this unit, you learned about cinematic motion and how to motivate it and do it properly. As such, you have gained understanding and practice with another important basic of excellent filmmaking.

\begin{enumerate}
\def\labelenumi{\arabic{enumi}.}
\tightlist
\item
  Types of Cinematic Motion and how to use tracking shots, following shots, dolly shots, amount others to add professionalism and more meaning in your projects.
\item
  Motivating Camera Movement where you explore the classic cinema technique of discussing your moving shots by having the camera and subject movement in sync.
\end{enumerate}

\hypertarget{assessment-7}{%
\section*{Assessment}\label{assessment-7}}
\addcontentsline{toc}{section}{Assessment}

\begin{assessment}
{Course Journal}

After completing this unit, including the learning activities, you are asked to make sure you are doing journal entries and when appropriate to share your responses with your facilitator and classmates when you meet.

Note that entries are expected after every unit. Your journal reflections are submitted at the end of the course as part of the Journal One: Personal Journal and self assessment.

\emph{See the Assessments section for more details on submitting your journal, as well as the grading criteria.}

{Motivated Camera Movement Exercise}

\emph{This activity is ungraded, but part of journal entry}

For this exercise on motion, you will shoot \textbf{5 different types} of motivated kinetic shots, e.g., a pan, tilt, zoom, following action shot, etc. that have smooth motion from the beginning to the end of the shot and ``disguise'' the camera technique by matching the rate of camera move and subject movement.

See the Assessments section for more details, as well as the grading criteria.
\end{assessment}

\hypertarget{checking-your-learning-3}{%
\section*{Checking your Learning}\label{checking-your-learning-3}}
\addcontentsline{toc}{section}{Checking your Learning}

\begin{progress}
Before you move on to the next unit, you may want to check to make sure that you are able to:

\begin{itemize}
\tightlist
\item
  Describe the various types of film motion shots
\item
  Contrast the various types of film motion shots
\item
  Determine when and how to use cinematic motion
\item
  Create effective cinematic motion shots
\end{itemize}
\end{progress}

\hypertarget{the-grammar-of-film}{%
\chapter{The Grammar of Film}\label{the-grammar-of-film}}

\hypertarget{overview-4}{%
\section*{Overview}\label{overview-4}}
\addcontentsline{toc}{section}{Overview}

We go back to theory in this unit and then have a chance to practice and apply it.

When we watch a film or TV show most people have no idea how much work and detail it took to make the production. Nor do they know the rules, guidelines, principles, and practices that were followed to create it. This unit and the ones following will help you explore what these best practices are and understand how to apply them to your work.

This will be a fun chapter as you get to explore and experiment with the fundamental building blocks that create meaning in cinema. It is a unit that will challenge both your thinking and your intuition skills as you learn the basic components of filmmaking and how to use them in creative ways.

\hypertarget{topics-4}{%
\subsection*{Topics}\label{topics-4}}
\addcontentsline{toc}{subsection}{Topics}

This unit is divided into the following topics:

\begin{enumerate}
\def\labelenumi{\arabic{enumi}.}
\tightlist
\item
  Basic Grammar of Film\\
\item
  Basic Shots and Their Descriptions\\
\item
  Camera Angles and Heights\\
\item
  Camera Subject Angles
\end{enumerate}

\hypertarget{learning-outcomes-4}{%
\subsection*{Learning Outcomes}\label{learning-outcomes-4}}
\addcontentsline{toc}{subsection}{Learning Outcomes}

When you have completed this unit, you should be able to:

\begin{itemize}
\tightlist
\item
  Define the basic elements of the grammar of film\\
\item
  Describe basic shots and their descriptions\\
\item
  Contrast camera angles and heights\\
\item
  Determine when and how to use subject angles\\
\item
  Create shots that reveal what you discovered in this unit
\end{itemize}

\hypertarget{activity-checklist-4}{%
\subsection*{Activity Checklist}\label{activity-checklist-4}}
\addcontentsline{toc}{subsection}{Activity Checklist}

\begin{reflect}
Here is a checklist of learning activities you will benefit from in completing this unit. You may find it useful for planning your work.

\begin{itemize}
\tightlist
\item
  Find and consult a film glossary such as Brooklyn College Film Glossary. Read and study the one you like throughout this course. Set a goal to learn 5 new terms or concepts each day.
\item
  Watch two videos about Basic Shots in action
\item
  Read and Study Chapter 5 and explain how frames and shots become scenes and sequences and how they add up to a complete film story.
\item
  After thoroughly reading and studying Chapter 5 to grasp the concepts of shots, camera angles, heights, and camera subject angles, engage in practical application.
\end{itemize}

\emph{Note that these learning activities are ungraded, unless specified, but are designed to help you succeed in your assessments in this course.}
\end{reflect}

\hypertarget{assessment-8}{%
\subsection*{Assessment}\label{assessment-8}}
\addcontentsline{toc}{subsection}{Assessment}

See the Assessment section in Moodle for assignment details and due dates.

\hypertarget{resources-4}{%
\subsection*{Resources}\label{resources-4}}
\addcontentsline{toc}{subsection}{Resources}

Here are the resources you will need to complete this unit.

\begin{itemize}
\tightlist
\item
  Chapter 5 of course text, \emph{Digital Filmmaking: A Beginner's Guide to Mastering the Craft.}
\item
  Other online resources will be provided in the unit.
\end{itemize}

\hypertarget{basic-grammar-of-film}{%
\section{Basic Grammar of Film}\label{basic-grammar-of-film}}

In the previous units we look at the basic building blocks of visual composition and cinematic motion. In this unit we are going to learn how to put them together with a variety of other shots to make cinematic context and meaning. This is the grammar of film used to make cinematic ``sentences'' which convey meaning and context to an audience as we create scenes, sequences, and short films.

However, as we focus on the components of filmmaking, we must not let the emphasis only be on film logic. We also want to cultivate the role that intuition and feeling play as we develop shots and put them together. This reminder is important because filmmaking is as much art as technical craft. We need to know the craft---the technical aspects---but also the art---the aesthetic and emotional dimensions. We need to keep this before us as we venture into the many elements and guidelines of effective filmmaking.

In this first topic you will learn about:
- Frames
- Shots
- Scenes
- Sequences
- Acts

These components give you the big picture of how the many smaller elements add up to create mood and meaning in a film project.

\hypertarget{activity-read-and-explore-1}{%
\subsection*{Activity: Read, and Explore}\label{activity-read-and-explore-1}}
\addcontentsline{toc}{subsection}{Activity: Read, and Explore}

\begin{reflect}
Find and consult a film glossary such as \href{http://userhome.brooklyn.cuny.edu/anthro/jbeatty/COURSES/glossary.htm}{Brooklyn College Film Glossary}. Continue to read and study the one you like throughout this course and it will help you grow your film vocabulary and your understanding of the many components and processes involved in filmmaking. Set a goal to learn 5 new terms or concepts each day.
\end{reflect}

\hypertarget{basic-shots-and-their-descriptions}{%
\section{Basic Shots and Their Descriptions}\label{basic-shots-and-their-descriptions}}

This is a vital section. You need to know the following shots and how to describe them because you will have to use this language to communicated to your crew. Few things are more frustrating for a film crew than the director or cinematographer not knowing what they want or how to name and describe a shot.

\begin{itemize}
\tightlist
\item
  Extreme Long Shot (ELS)\\
\item
  Long Shot (LS) also called Wide Shot (WS)\\
\item
  Establishing Shot (ES)\\
\item
  Full Shot (FS)\\
\item
  Medium Long Shot (MLS) also known as Medium Wide Shot (MWS)\\
\item
  Cowboy Shot\\
\item
  Medium Shot (MS)\\
\item
  Medium Close Up (MCU)\\
\item
  Close Up (CU)\\
\item
  Extreme Close Up (ECU)
\end{itemize}

In addition to shot size we also designate film shots by who or what is in the frame and how they are positioned:

\begin{itemize}
\tightlist
\item
  Single Shot\\
\item
  Two Shot\\
\item
  Three Shot\\
\item
  Group Shot\\
\item
  Over-The-Shoulder Shot\\
\item
  Reverse Shot\\
\item
  Point of View Shot (POV)\\
\item
  Reaction Shot\\
\item
  Cutaway\\
\item
  Cut-in and Insert\\
\item
  Combo Shot
\end{itemize}

\hypertarget{activity-watch-and-learn-1}{%
\subsection*{Activity: Watch and Learn}\label{activity-watch-and-learn-1}}
\addcontentsline{toc}{subsection}{Activity: Watch and Learn}

\begin{reflect}
As they say, a picture is worth a thousand words, and when you see the shots in action in these resources they become more clear. Watch the video resources below to learn more about basic shots

\href{https://www.youtube.com/watch?v=lRo2IqYbEaE}{Watch: Introduction to shot types and camera movement}

\href{https://www.youtube.com/watch?v=7y0ouVBcogU}{Watch: 15 Essential Camera Shots, Angles and Movements in Filmmaking}
\end{reflect}

\hypertarget{camera-angles-and-heights}{%
\section{Camera, Angles, and Heights}\label{camera-angles-and-heights}}

Camera heights and angles reveal both text and subtext. That is, they show us what is being filmed (the object or subject) and often convey something about the meaning of the subject or object (the subtext).

In essence, there are no ``neutral'' camera angles. Our camera height placement will convey something and it is important for the filmmaker to know what that something is.

The different camera heights include:

\begin{itemize}
\tightlist
\item
  Eye-Level\\
\item
  Low Angle\\
\item
  High Angle\\
\item
  Shoulder Level\\
\item
  Hip Level\\
\item
  Knee Level Shot\\
\item
  Ground Level\\
\item
  Worm's Eye View\\
\item
  Bird's Eye View\\
\item
  God's Eye View\\
\item
  Dutch Tilt or Dutch Canted Angle
\end{itemize}

\hypertarget{activity-read-and-study-1}{%
\subsection*{Activity: Read and Study}\label{activity-read-and-study-1}}
\addcontentsline{toc}{subsection}{Activity: Read and Study}

\begin{reflect}
Read and study Chapter 5. In your Film Journal, explain how frames and shots become scenes and sequences and how they add up to a complete film story. Also, summarize the best thing you learned about camera heights in this section.

Commit to doing one technique several times a day, and reflect on what you learnt from this experience.

\textbf{Helpful Hint:} It is recommended but not mandatory that you study a web essay on how to build your intuition such as\href{https://liveboldandbloom.com/04/self-improvement/develop-your-intuition}{21 Eye-Opening Ways To Develop Your Intuition}
\end{reflect}

\hypertarget{camera-subject-angles}{%
\section{Camera Subject Angles}\label{camera-subject-angles}}

Subject angle refers to the perspective the audience has of the person being filmed. Subject angles are important because they affect the way the audience interprets the character or scene.

The following are key subject angles:

\begin{itemize}
\tightlist
\item
  Profile\\
\item
  Full-Face\\
\item
  Three-quarter\\
\item
  Quarter angle\\
\item
  Full-Back\\
\item
  Combo\\
\item
  Silhouettes
\end{itemize}

Having completed this topic, you now have a firm grasp of the grammar of film and its components.

\hypertarget{activity-practice}{%
\subsection*{Activity: Practice}\label{activity-practice}}
\addcontentsline{toc}{subsection}{Activity: Practice}

\begin{reflect}
After thoroughly reading and studying Chapter 5, it's time to put theory into practice. Engage in practical application using camera angles, heights and camera subjects
\end{reflect}

\hypertarget{summary-4}{%
\section*{Summary}\label{summary-4}}
\addcontentsline{toc}{section}{Summary}

In this unit, you learned about the importance of film grammar and the multitude of shot sizes, framings, angles, etc. In doing so, you now have a good foundational grasp of the basic components needed to make excellent film projects.

\hypertarget{assessment-9}{%
\section*{Assessment}\label{assessment-9}}
\addcontentsline{toc}{section}{Assessment}

\begin{assessment}
{Course Journal}

After completing this unit, including the learning activities, you are asked to make sure you are doing journal entries and when appropriate to share your responses with your facilitator and classmates when you meet.

Note that entries are expected after every unit. Your journal reflections are submitted at the end of the course as part of the Journal One: Personal Journal and self assessment.

\emph{See the Assessments section for more details on submitting your journal, as well as the grading criteria.}
\end{assessment}

\hypertarget{checking-your-learning-4}{%
\section*{Checking your Learning}\label{checking-your-learning-4}}
\addcontentsline{toc}{section}{Checking your Learning}

\begin{progress}
Before you move on to the next unit, you may want to check to make sure that you are able to:

\begin{itemize}
\tightlist
\item
  Define the basic elements of the grammar of film.\\
\item
  Describe basic shots and their descriptions.\\
\item
  Contrast camera angles and heights.\\
\item
  Determine when and how to use subject angles.\\
\item
  Create shots that reveal what you discovered in this unit.
\end{itemize}
\end{progress}

\hypertarget{visual-storytelling}{%
\chapter{Visual Storytelling}\label{visual-storytelling}}

\hypertarget{overview-5}{%
\section*{Overview}\label{overview-5}}
\addcontentsline{toc}{section}{Overview}

Congratulations. You are now ready to make a rare type of film. Rare you say? Yes, because you are going to shoot your film project in a way that is not done very often, that is, in full sequential order. What this means will be explained ahead.

We are now at a point in the course where you are going to be held accountable for your work. The earlier exercises were not graded to give you the chance to explore and make mistakes, without regard to a grade. We now have to cross into that territory. The main reason for this is that, as mentioned previously, film is a public medium and your work will be critiqued when it is shown. So this will be a great opportunity to learn to give and receive feedback.

But fear not. Based on your previous work and what you will learn in this unit you will be ready to make a film that is ``public'' worthy. In addition, this first film will only account for 10\% of your grade so you are encouraged to explore and take risks.

\hypertarget{topics-5}{%
\subsection*{Topics}\label{topics-5}}
\addcontentsline{toc}{subsection}{Topics}

This unit is divided into the following topics:

\begin{itemize}
\tightlist
\item
  In-Camera Project and Its Benefits
\item
  Brainstorming and Successful Creativity
\item
  Direction Vectors and Eye-line Requirements
\end{itemize}

\hypertarget{learning-outcomes-5}{%
\subsection*{Learning Outcomes}\label{learning-outcomes-5}}
\addcontentsline{toc}{subsection}{Learning Outcomes}

When you have completed this unit, you should be able to:

\begin{itemize}
\tightlist
\item
  Describe what constitutes a liner story\\
\item
  Define what constitutes an ``In-Camera'' film\\
\item
  Analyze and apply a story template\\
\item
  Determine how to tell a visually-centered story well\\
\item
  Create a short simple story to film
\end{itemize}

\hypertarget{activity-checklist-5}{%
\subsection*{Activity Checklist}\label{activity-checklist-5}}
\addcontentsline{toc}{subsection}{Activity Checklist}

\begin{reflect}
Here is a checklist of learning activities you will benefit from in completing this unit. You may find it useful for planning your work.

\begin{itemize}
\tightlist
\item
  Read, Study, and Reflect: Read and study Chapter 6 of course e-text: \emph{Digital Filmmaking: A Beginner's Guide to Mastering the Craft.} Write in your journal your initial thoughts regarding the benefits of doing an in-camera film
\item
  Brainstorming: Brainstorm at least five story ideas and use the ``Fairy Tale Template for a Short Film'' to write your story for your short film.
\end{itemize}

\emph{Note that these learning activities are ungraded (unless specified), but designed to help you succeed in your assessments in this course}
\end{reflect}

\hypertarget{assessment-10}{%
\subsection*{Assessment}\label{assessment-10}}
\addcontentsline{toc}{subsection}{Assessment}

See the Assessment section in Moodle for assignment details and due dates.

\hypertarget{resources-5}{%
\subsection*{Resources}\label{resources-5}}
\addcontentsline{toc}{subsection}{Resources}

Here are the resources you will need to complete this unit.

\begin{itemize}
\tightlist
\item
  Chapter 6 of course text: \emph{Digital Filmmaking: A Beginner's Guide to Mastering the Craft,} by Ned Vankevich (e-text)
\item
  Other online resources will be provided in the unit.
\end{itemize}

\hypertarget{in-camera-project-and-its-benefits}{%
\section{In-Camera Project and Its Benefits}\label{in-camera-project-and-its-benefits}}

Your first film project will be something you might never do again if you go into the film and television world. You will shoot a short film with no dialogue in the exact sequential order that it will appear when you screen it for classs.

Almost always, films are shot out of sequence to save time and money by doing all the scenes in one location at a time and then all the scenes in another location at one time, etc.

Here you will be forced to shoot at a location and move to the next, and if the first location is needed you will have to go back to it. Why do this? There are multiple benefits, as the course text chapter highlights. Some of them include:

\begin{itemize}
\tightlist
\item
  Forcing you to think of a simple story that fits these parameters.\\
\item
  Challenging you to plan your shots in a linear, well-thought out way.\\
\item
  Making you attentive to each shot and how relates to the previos ones.\\
\item
  Challenging your brainstorming and film logic skills.
\end{itemize}

\hypertarget{activity-read-study-and-reflect}{%
\subsection*{Activity: Read, Study and Reflect}\label{activity-read-study-and-reflect}}
\addcontentsline{toc}{subsection}{Activity: Read, Study and Reflect}

\begin{reflect}
Read and study Chapter 6. Write in your journal your initial thoughts regarding the benefits of doing an in-camera film. Does it thrill and excite you or make you apprehensive and anxious? Or maybe a combination of these feelings and emotions. As artists we need to get in touch with our feelings and learn to work with them not fear them.

To prime your imagination for the upcoming project, watch the video \href{https://www.youtube.com/watch?v=iWQQgZh9EyE}{Visual Storytelling 101}
\end{reflect}

{ Secrets to a Simple Story}

In order to do this first project, you will need to come up with a story to film. Something has to happened to someone and they must engage in action to address it. But if we are to design a good story---one that holds the attention of our audience and engages them---we need to focus on several things in our short film project:

\begin{itemize}
\tightlist
\item
  One or two main characters.\\
\item
  A single problem or conflict.\\
\item
  Knowing your genre (will your film be funny, dramatic, scary, etc.?).\\
\item
  A simple setting or number of locations.\\
\item
  A satisfying ending.
\end{itemize}

You need to keep this simple formula in mind in the projects ahead. There will be many forces and temptations that will distract you and pull you away from this simple understanding. Resist them and go back to basics if you get lost, confused, or frustrated.

\hypertarget{brainstorming-and-successful-creativity}{%
\section{Brainstorming and Successful Creativity}\label{brainstorming-and-successful-creativity}}

The secret to a story sounds simple and it is. We need to create a character with a goal and stakes (the painful consequences that will happen to the character) who has to overcome opposition to that goal and an ending that reveal whether or not the character obtains the goal.

However, the challenge for most of us is, what story do I tell?

Brainstorming is a great tool and way to discover and find the right story you want to tell. Brainstorming involves the spontaneous development of ideas. Brainstorming can be done alone or within a group.

The great thing about brainstorming is that you do not have to judge and criticize the process and results. In addition, you are not trying to perfect. You are only looking for that one idea that sparks your creativity and you ignore the others. This non-judgmentalism and jettisoning of perfectionism is important because these two tendencies stifle creativity.

Even if you already have a story idea that you want to film you are strongly encouraged to engage in brainstorming activities during this section and the rest of the course. If you do, you unleash more of your creativity and you will most likely find a better story concept.

\hypertarget{activity-brainstorming}{%
\subsection*{Activity: Brainstorming}\label{activity-brainstorming}}
\addcontentsline{toc}{subsection}{Activity: Brainstorming}

\begin{reflect}
After you feel you have a good grasp of brainstorming, engage in the process and come up with 5 story concepts that fit the criteria in the previous topic:

\begin{itemize}
\tightlist
\item
  One or two main characters.
\item
  A single problem or conflict.
\item
  Knowing your genre (will your film be funny, dramatic, scary, etc.?).
\item
  A simple setting or number of locations.
\item
  A satisfying ending. After you have done this go through your concept list and choose your top one and apply it to the template in the next topic.
\end{itemize}

\textbf{Helpful Hint:} If you find yourself stuck and can't come up with a story concept, tell your facilitator and have a group brainstorming session. Take one of the ideas and develop it. Share your concept if you are stuck finding one of the criteria such as a goal or stakes and ask for brainstorming help to solve the problem. Remember, film is mostly a collaborative venture.
\end{reflect}

{Short Film Story Template}

Once you have your core story concept you now need to make it into a full story with a beginning, middle, and end. We also need to make sure we do not make it too complicated or muddled.

Many of us have the tendency to overly complicate things: to add extraneous detail or to have elements that are not clear well explained. The template is this section is designed to prevent these issues as you use a fairy tale structure to find your characters, the problem they are facing, and show what happens in a clear way with a beginning, middle, and end of the story.

Have fun with this template process. Also do not forget to use brainstorming for each section of the template so that you come up with the best characters, situations, locations, and conflict which are not predicable or have a ``been there, done that'' (boring) quality.

\hypertarget{activity-planning-your-film}{%
\subsection*{Activity: Planning Your Film}\label{activity-planning-your-film}}
\addcontentsline{toc}{subsection}{Activity: Planning Your Film}

\begin{reflect}
Once you have brainstormed and have a clear and solid concept film, e.g., a student who has the power to turn things in gold, use the \href{assets/unit6/Fairy_Tale_Template_for_a_Short_Film.pdf}{Fairy Tale Template for a Short Film} and write your story for your short film by filling in the blanks.

Study the illustration example given in the chapter and then apply the process to the creation of your film story. In crafting your story you are not allowed to use dialogue beyond ``yes,'' ``no,'' and ``okay.'' That is, you must find visual ways to establish your story's setup and context and its build and payoff. For example, you cannot have your character say ``I am on my way to the store.'' You will just show him or her getting in the car, driving, arriving at the store parking lot, and entering the store. (This has visual direction vectors discussed below in Topic 3.)
\end{reflect}

\hypertarget{direction-vectors-and-eyeline-requirements}{%
\section{Direction Vectors and Eyeline Requirements}\label{direction-vectors-and-eyeline-requirements}}

Now that you have your story and one-line action descriptions, and have your actors, props, and locations set-up, you are ready to film your story (the In-Camera film project) with your cell phone.

Before doing so, this exercise will require that you focus on several things found in the Direction Vectors and Eyeline Requirements section:

\begin{itemize}
\tightlist
\item
  Proper Lead Room\\
\item
  Motivated Shot Movement\\
\item
  Consistent Screen Direction\\
\item
  Proper Eye-Line Vectors\\
\item
  Proper Headroom
\end{itemize}

These requirements should not worry you since you have explored and tried many of these techniques in the composition and film movements units.

\hypertarget{activity-resources-on-direction-vectors-and-eyeline-requirements}{%
\subsection*{Activity: Resources on Direction Vectors and Eyeline Requirements}\label{activity-resources-on-direction-vectors-and-eyeline-requirements}}
\addcontentsline{toc}{subsection}{Activity: Resources on Direction Vectors and Eyeline Requirements}

\begin{reflect}
In order to do well and follow the direction vectors and eyeline requirements for this first film project, be sure to consult the following resources which explain them in visual ways. As you do so, pre-visualize how you will film your project:

\href{https://www.youtube.com/watch?v=y_1H6V7uyYc}{Watch: The eyeline match}

\begin{center}\rule{0.5\linewidth}{0.5pt}\end{center}

\href{https://www.youtube.com/watch?v=9XOn5uxdSJc}{Watch: Screen Direction rule}

\begin{center}\rule{0.5\linewidth}{0.5pt}\end{center}

In addition, be sure to review Units 3, 4, and 5 to make sure your shots will be well composed and that you will incorporate a variety of camera angles.

Happy filming!
\end{reflect}

\hypertarget{summary-5}{%
\section*{Summary}\label{summary-5}}
\addcontentsline{toc}{section}{Summary}

In this unit, you learned about:

\begin{itemize}
\item
  The In-Camera Project and how It is beneficial to help you to understand how to think of filmmaking in a linear way.
\item
  Secrets to a Simple Story and how limiting your story to a key single concept or theme with only a few characters and locations will help you create a stronger film.
\item
  Brainstorming and Successful Creativity and how the more ideas and concepts you imagine and generate, the better your film will capture an audience's attention.
\item
  Short Story Film Template where a fairy tale story formula can help you design the beginning, middle, and end of your film.
\item
  Direction Vectors and Eye-line Requirements which are important cinematic compositional devices that prevent the audience's confusion as you transition between shots.
\end{itemize}

\hypertarget{assessment-11}{%
\section*{Assessment}\label{assessment-11}}
\addcontentsline{toc}{section}{Assessment}

\begin{assessment}
{Course Journal}

After completing this unit, including the learning activities, you are asked to make sure you are doing journal entries and when appropriate to share your responses with your facilitator and classmates when you meet.

Also, log in your journal what you learned from the creating and filming of your In-Camera exercise. Log what you thought about the films of your peers and what you learned from the feedback from them and your instructor.

Be sure to make a note of what films stood out and who did them. You will consult this at the end of the course when the class determines: Best Overall Film. Most Imaginative Film, Best Story, Best Cinematography, Best Editing, etc.

\emph{See the Assessments section for more details on submitting your journal, as well as the grading criteria.}

\begin{center}\rule{0.5\linewidth}{0.5pt}\end{center}

{In-Camera Exercise}

\emph{This activity is ungraded, but important}

Create a brief mini movie that shows your understanding of how to apply composition, motivated camera movement, and a variety of shots to visually tell a story. That is, images and not dialogue should drive your story.

See the Assessments section for more details on submitting your journal, as well as the grading criteria.
\end{assessment}

\hypertarget{checking-your-learning-5}{%
\section*{Checking your Learning}\label{checking-your-learning-5}}
\addcontentsline{toc}{section}{Checking your Learning}

\begin{progress}
Before you move on to the next unit, you may want to check to make sure that you are able to:

\begin{itemize}
\tightlist
\item
  Describe what constitutes a good, basic, and simple story\\
\item
  Define what constitutes an ``In-Camera'' film\\
\item
  Analyze and apply a story template\\
\item
  Determine how to tell a visual-centered story well\\
\item
  Create a short simple story to film
\end{itemize}
\end{progress}

\hypertarget{editing-and-montage}{%
\chapter{Editing and Montage}\label{editing-and-montage}}

\hypertarget{overview-6}{%
\section*{Overview}\label{overview-6}}
\addcontentsline{toc}{section}{Overview}

Congratulations. You've made your first mini-movie in this course. Don't worry if it did not live up to your expectations. The important thing right now is that you wrote, produced, directed, and shot a short film and that you are improving your skills. The rest of your course will give you a chance to explore more elements of filmmaking and gain more experience.

In this unit you will focus on editing and creating a montage, a common form of filmmaking.

Unlike the previous project, you can now use dialogue, voice over, music, and sound effects to enhance your short film project. You will also use editing which will allow you to shoot your film out of order and then assemble it after you have all of the elements you need to make it.

As you saw in Unit 2, editing is the third major phase of filmmaking and a major skill to understand.

Don't worry if you do not want to be an editor. You will work in teams this time and as long as the work on the film is shared equally between yourself and your partner(s), you can offer editing suggestions but will not have to push the buttons. (This is how it works in the industry where the director and producer tell the editor what they want and the editor does the technical work.)

\hypertarget{topics-6}{%
\subsection*{Topics}\label{topics-6}}
\addcontentsline{toc}{subsection}{Topics}

This unit is divided into the following topics covered in the course text:

\begin{itemize}
\tightlist
\item
  Post-Production Workflow
\item
  Montage
\item
  Digital Video Editing Software
\item
  Montage Project
\end{itemize}

\hypertarget{learning-outcomes-6}{%
\subsection*{Learning Outcomes}\label{learning-outcomes-6}}
\addcontentsline{toc}{subsection}{Learning Outcomes}

When you have completed this unit, you should be able to:

\begin{itemize}
\tightlist
\item
  Define the Post-Production workflow
\item
  Analyze the necessary elements needed for your montage
\item
  Determine the steps needed in making a montage
\item
  Create an effective montage
\item
  Evaluate the quality of a montage
\end{itemize}

\hypertarget{activity-checklist-6}{%
\subsection*{Activity Checklist}\label{activity-checklist-6}}
\addcontentsline{toc}{subsection}{Activity Checklist}

\begin{reflect}
Here is a checklist of learning activities you will benefit from in completing this unit. You may find it useful for planning your work.

\begin{itemize}
\item
  Read and study Chapter7 of course etext: \emph{Digital Filmmaking: A Beginner's Guide to Mastering the Craft.} The readings are divided into 6 topics:

  \begin{itemize}
  \tightlist
  \item
    The Post-Production Workflow
  \item
    Montage
  \item
    Types of Montages
  \item
    Picture Cutting Techniques.
  \item
    Digital Editing Software
  \item
    Creating a Montage
  \end{itemize}
\item
  Watch videos related to the topics previously mentioned.
\item
  As you read Chapter 7 ``Digital Editing Software'' section, find the digital video editing program you want to use and learn how to do basic editing with it.
\item
  Practice doing your pitch before a friend or family member or to an imaginary audience before making the actual pitch to your course members.
\end{itemize}
\end{reflect}

\hypertarget{assessment-12}{%
\subsection*{Assessment}\label{assessment-12}}
\addcontentsline{toc}{subsection}{Assessment}

See the Assessment section in Moodle for assignment details and due dates.

\hypertarget{resources-6}{%
\subsection*{Resources}\label{resources-6}}
\addcontentsline{toc}{subsection}{Resources}

Here are the resources you will need to complete this unit.

\begin{itemize}
\tightlist
\item
  Chapter 7 of course text: \emph{Digital Filmmaking: A Beginner's Guide to Mastering the Craft,} by Ned Vankevich (e-text)
\item
  Other online resources will be provided in the unit.
\end{itemize}

\hypertarget{post-production-workflow}{%
\section{Post-Production Workflow}\label{post-production-workflow}}

There are many steps and elements involved in the post-production process. They include:

\begin{itemize}
\tightlist
\item
  Importing Footage and elements\\
\item
  Sound syncing\\
\item
  Picture cutting\\
\item
  Transitions\\
\item
  Dialogue editing\\
\item
  Sound design\\
\item
  Adding music\\
\item
  Adding sound effects\\
\item
  Adding visual and special effects\\
\item
  Color correcting\\
\item
  Adding tiles and graphics\\
\item
  Etc.
\end{itemize}

\hypertarget{activity-read-and-watch}{%
\subsection*{Activity: Read and Watch}\label{activity-read-and-watch}}
\addcontentsline{toc}{subsection}{Activity: Read and Watch}

\begin{reflect}
Before completing your reading for this unit, log in your journal what type of film you want to work on for your final film project and why.

Then, read the first section in Chapter 7, the Post-Production Workflow. For a detailed overview of the post-production process, watch the video \href{https://www.youtube.com/watch?v=IFjyVW21Vpw}{''Stages of Post-Production for Filmmaking in Cinema''}
\end{reflect}

\hypertarget{montage}{%
\section{Montage}\label{montage}}

Montage involves a type of editing where a series of images and sound elements are most often used used to condense time, create emotion, tell a story, reveal something from the past, promote something, or impart information.

Montages can stand alone or be part of a larger documentary, film or television story, or a stand-alone film.

You will focus on a montage project in this unit because they are a great way to learn the art of editing when you shoot images and edit them in timed rhythm to music, dialogue, and/or sound effects. In this exercise you are going to let emotions, feelings, and mood guide you as you learn to further develop your intuitive creative skills.

\hypertarget{picture-cutting-techniques}{%
\section{Picture Cutting Techniques}\label{picture-cutting-techniques}}

This montage project is helpful because you have to focus on a few elements: the images you will use and the music and/or sound effects that will provide the beats and rhythm for your picture cutting.

Although the music will guide your editing as you cut on beats, the images you edit will be important. Once again, good composition and motivate motions shots will play a vital role as you shot your footage for your montage project. You can also used ``found'' or archival footage, or news stories as your visual sources in a montage.

There are many types of cuts and transitions that can be used in a montage as you splice and stich your footage together including:

\begin{itemize}
\tightlist
\item
  Cuts
\item
  Jump-Cuts
\item
  Cross-Cutting
\item
  Dissolves
\item
  Wipes
\item
  Fades
\end{itemize}

You can use J-cut and L-cuts when you are working with dialogue or you want to lead into a new shot or scene or carry the audio to a new shot or scene. If you are interested in doing this, watch the video \href{https://www.youtube.com/watch?v=PAvJevWUVsc}{``Film Making 101 Tutorial:L\&J Cuts''}

\hypertarget{types-of-montage-projects}{%
\section{Types of Montage Projects}\label{types-of-montage-projects}}

You are given creative latitude and freedom for the type of montage you will shoot and edit for this unit. Some of the most popular ones include:

\begin{itemize}
\tightlist
\item
  Poetic Montages
\item
  Nature Montages
\item
  Street Montages
\item
  Sequence Sprint
\item
  Music Videos
\item
  Aesthetic Montages
\end{itemize}

You can also combine these types of montages. As you read the description of these different types, have fun choosing the one you want to do.

\hypertarget{activity-read-and-study-2}{%
\subsection*{Activity: Read and Study}\label{activity-read-and-study-2}}
\addcontentsline{toc}{subsection}{Activity: Read and Study}

\begin{reflect}
Read and study the second section in Chapter 7, Montage, types of Montages, and Picture Cutting Techniques.
\end{reflect}

\hypertarget{activity-watch-and-journal}{%
\subsection*{Activity: Watch and Journal}\label{activity-watch-and-journal}}
\addcontentsline{toc}{subsection}{Activity: Watch and Journal}

\begin{reflect}
For visual examples of Picture Cutting Techniques, watch the videos \href{https://www.youtube.com/watch?v=OAH0MoAv2CI}{Cuts \& Transitions 101} and \href{https://www.youtube.com/watch?v=Wv3Hmf2Dxlo}{9 Cuts Every Video Editor Should Know}

\begin{center}\rule{0.5\linewidth}{0.5pt}\end{center}

\begin{center}\rule{0.5\linewidth}{0.5pt}\end{center}

\textbf{Optional}: For an excellent example of a nature and time-lapse montage where a series of images are synced into a powerful soundtrack , watch the video \href{https://www.youtube.com/watch?v=oagszCmJLpU}{``Pursuit - A 4K storm time-lapse film''.} As you view it, pay attention to the way the types and rhythms of the images flow and work in sync with the music and sound effects especially from 2: 29 onward.

Lastly, log in in your journal what you learn about this montage
\end{reflect}

\hypertarget{digital-video-editing-software}{%
\section{Digital Video Editing Software}\label{digital-video-editing-software}}

You are now at a point where you will have to import your cellphone or DSLR footage into a software program that will allow you to edit it. Digital video editing software is helpful because it allows you to edit your images and sync them to the rhythm and beats of your sound track and export your project for viewing or streaming on the Internet.

There are many types and brands of editing software in the marketplace and you will have to choose the one that works for you. Some applications only work with Android-based cellphones and others for Apple ISO devices and some for both platforms such as imovie. If you want to really get series about editing you should explore Final Cut Pro, Adobe Premiere, Avid, DaVinci Resolve, or some other professional program.

There are similarities among editing software and programs but you will have to do a tutorial on YouTube or the manufacture's website of the one you will use to see how to work it. The more popular the software, the more tutorials will be available.

\hypertarget{activity-read-study-and-explore}{%
\subsection*{Activity: Read, Study, and Explore}\label{activity-read-study-and-explore}}
\addcontentsline{toc}{subsection}{Activity: Read, Study, and Explore}

\begin{reflect}
Read and study Chapter 7, Digital Editing Software section. Next, Find the digital video editing program you want to use and learn how to do basic editing with it. Basic editing using cuts, dissolves, fades, and wipes titling will be all you need to learn for your montage project.

The \href{https://filmora.wondershare.net/filmora-video-editor.html?}{Filmora} website gives an overview of the type of features you want to look for. For free and inexpensive Android-based video software see:\href{https://fossbytes.com/best-android-video-editor-apps/}{13 Best Android Video Editor Apps Of 2022} or search for best cellphone or DSLR video editing software such as \href{https://www.tomsguide.com/best-picks/best-video-editing-apps}{Best video editing apps 2023}
\end{reflect}

\hypertarget{montage-project}{%
\section{Montage Project}\label{montage-project}}

{Creating Your Montage}

As we have seen, there are three main phases of film production and they will apply here:

\begin{itemize}
\tightlist
\item
  Montage Pre-Production
\item
  Montage Production
\item
  Montage Post-Production
\end{itemize}

The lion's share of the emphasis in the chapter has been on post-production, but this does not mean you must take your pre-production and production work lightly. Your montage will only be as good as the images you have for it.

The following criteria can help us better evaluate our own work and that of others and to give constructive criticism on how to improve:

\begin{itemize}
\tightlist
\item
  Overall did the montage work?
\item
  Did the music and/or sound effects work well with the images?
\item
  Did the montage convey a mood, emotion, a story and/or theme (some insight into life)? If so, what was it?
\item
  Were the shots well composed, motivated, and appropriately smooth or in some cases jarring depending on what is being communicated?
\item
  Were the cuts timed well? If not, where did they not work well?
\item
  Were the transitions appropriate and effective?
\item
  Was the action followed well?
\item
  Was there a sufficient variety of shots?
\item
  Were some shots or images too repetitive?
\item
  Were the shots and images interesting and noteworthy?
\item
  What did you like best about the project?
\item
  What suggestions would you offer to make it better?
\end{itemize}

\hypertarget{activity-read-and-study-3}{%
\subsection*{Activity: Read and Study}\label{activity-read-and-study-3}}
\addcontentsline{toc}{subsection}{Activity: Read and Study}

\begin{reflect}
Study the next section of Chapter 7 on creating a montage. As you study montage pre-production, production, and post-production make notes on your shooting script of what to watch out for and not forget as you film the shots and images you will use.
\end{reflect}

\hypertarget{activity-creating-a-montage}{%
\subsection*{Activity: Creating a Montage}\label{activity-creating-a-montage}}
\addcontentsline{toc}{subsection}{Activity: Creating a Montage}

\begin{reflect}
Once you have edited your project and exported it, you are ready to show it to your classmates and instructor. This should be an exciting time as you screen the fruits of your labor and creativity! As previously mentioned, film is a public medium and subject to critique and evaluation.

Evaluate and critique the work of your peers helping them to see what worked and did not work in their montage and why. When possible give suggestions of how something could have been done better. Remember the tone and substance of our critiques should be one that offers encouraging constructive insight and that helps each other to grow and improve their craft.

For insight into how to critique a short film watch the video \href{https://www.youtube.com/watch?v=x5ASDkOSIBE}{``Critiquing Your Short Films''}
\end{reflect}

\hypertarget{activity-self-assessment}{%
\subsection*{Activity: Self-Assessment}\label{activity-self-assessment}}
\addcontentsline{toc}{subsection}{Activity: Self-Assessment}

\begin{reflect}
Write in your journal what you learned about your project and yourself from the montage projects.

Use the following questions to guide you:

\begin{itemize}
\tightlist
\item
  What was the best lesson I learned?\\
\item
  What did I do well? Why?\\
\item
  What did not work out as well as I planned? Why?\\
\item
  What can I do better next time?\\
\item
  What was the quality of my experience working with a partner?\\
\item
  If I did not work with a partner, could the project have been better if I had one?\\
\item
  How did this experience help me grow as a person and as a professional?\\
\item
  What did I learn new about myself?
\end{itemize}

Be sure to note which films and filmmakers had the best cinematography, editing, story, creativity, etc. so that you can consult this when you vote on the awards at the end of the course.
\end{reflect}

\hypertarget{summary-6}{%
\section*{Summary}\label{summary-6}}
\addcontentsline{toc}{section}{Summary}

In this unit, you learned about:

\begin{itemize}
\tightlist
\item
  Post-Production Workflow which helps you to organize the process of how you will edit your film project.
\item
  Montage which is a technique or type of film where pictures and sounds are interwoven to create a story, mood, and theme.
\item
  Picture Cutting Techniques that are the transitions you will need to move from shot to shot and scene to scene.
\item
  Types of Montage include sports action, fight and chase scenes, parallel stories, film poems, among others.
\item
  Digital Video Editing Software that you must find that is appropriate for your computer, tablet, or cell phone operating system.
\item
  Steps for Creating Your Montage where you develop your story, plan and shoot your shots, find the appropriate music and sound effects and create the appropriate transition for them.
\item
  Pitching Your Montage Project where you explain your story idea and why it is worth doing and how you will develop and produce it.
\item
  Feedback and Self-Assessment for Montage Project where your project will be critiqued by your professor and your peers; and you will explain what you learned from this experience.
\end{itemize}

\hypertarget{assessment-13}{%
\section*{Assessment}\label{assessment-13}}
\addcontentsline{toc}{section}{Assessment}

\begin{assessment}
{Course Journal}

After completing this unit, including the learning activities, you are asked to make sure you are doing journal entries and when appropriate to share your responses with your facilitator and classmates when you meet.

Note that entries are expected after every unit. Your journal reflections are submitted at the end of the course as part of the Journal One: Personal Journal and self assessment.

\emph{See the Assessments section for more details on submitting your journal, as well as the grading criteria.}

\begin{center}\rule{0.5\linewidth}{0.5pt}\end{center}

{Montage Short Project (30\%)}

This montage project will account for 30 \% of your grade and will be broken down according to the following:

\begin{itemize}
\tightlist
\item
  Story/Montage Concept: 20\%\\
\item
  Cinematography: 20\%\\
\item
  Editing: 40\%\\
\item
  Overall all project quality: 20\%
\end{itemize}

See the Assessment section of the course for more details.
\end{assessment}

\hypertarget{checking-your-learning-6}{%
\section*{Checking your Learning}\label{checking-your-learning-6}}
\addcontentsline{toc}{section}{Checking your Learning}

\begin{progress}
Before you move on to the next unit, you may want to check to make sure that you are able to:

\begin{itemize}
\tightlist
\item
  Define the Post-Production workflow\\
\item
  Describe how to pitch a project well\\
\item
  Analyze the necessary elements needed for your montage\\
\item
  Determine the steps needed in making a montage\\
\item
  Create an effective montage\\
\item
  Evaluate the quality of a montage
\end{itemize}
\end{progress}

\hypertarget{creating-narrative-and-other-short-films}{%
\chapter{Creating Narrative and Other Short Films}\label{creating-narrative-and-other-short-films}}

\hypertarget{overview-7}{%
\section*{Overview}\label{overview-7}}
\addcontentsline{toc}{section}{Overview}

More kudos to you. You now have two film projects and a lot of film exercises in your experience bank account. You will now be moving up a level where you will have the opportunity to bring together all of what you have learned so far into a longer narrative film.

In the previous units you focused on how to make shorter films. In this unit and the next you will focus on creating longer forms where your film will be five to ten minutes in length. This might not seem long but the longer the film, the more you have to do to sustain interest and continually engage the audience. Mini-films (2-minutes and under) have less of this burden because an audience is not feeling like they are wasting a bigger chunk of their time if the film is not well done.

Narrative films, whether fiction or non-fiction, are popular because they tell stories. Even if you do not want to be a ``Hollywood'' filmmaker learning to tell film stories will aid you in your personal and professional development. Human beings are hard-wired for storytelling. It is a primary way we make sense of the world and pass our knowledge and cultural values and traditions to others. Those who control the narratives of a society have powerful influence and impact.

Before moving ahead, view this resource:

\href{https://www.youtube.com/watch?v=mYnsKATCrdw}{7 Things to Know About Making Short Films:FRIDAY 101}

\hypertarget{topics-7}{%
\subsection*{Topics}\label{topics-7}}
\addcontentsline{toc}{subsection}{Topics}

This unit is divided into the following topics:

\begin{itemize}
\tightlist
\item
  Types of films and genres
\item
  The Logistics for the Final Project\\
\item
  Script and Story Development\\
\item
  Creating a Scriptment
\end{itemize}

\hypertarget{learning-outcomes-7}{%
\subsection*{Learning Outcomes}\label{learning-outcomes-7}}
\addcontentsline{toc}{subsection}{Learning Outcomes}

When you have completed this unit, you should be able to:

\begin{itemize}
\tightlist
\item
  Describe the different types of short films\\
\item
  Define film genre and why it is important\\
\item
  Determine how to approach making a film script to shoot and edit.\\
\item
  Create a shooting script
\end{itemize}

\hypertarget{activity-checklist-7}{%
\subsection*{Activity Checklist}\label{activity-checklist-7}}
\addcontentsline{toc}{subsection}{Activity Checklist}

\begin{reflect}
Here is a checklist of learning activities you will benefit from in completing this unit. You may find it useful for planning your work.

\begin{itemize}
\tightlist
\item
  Log in your journal what type of film you want to work on for your final film project and why.
\item
  Review popular short film genres.
\item
  Follow the steps in \emph{Story Research} section of Chapter Eight and ask and answer the questions in the \emph{What to Look for} section. View the video selected.
\item
  Brainstorm and develop a core story idea to pitch to the class.
\item
  Create a scriptment with your group.
\end{itemize}
\end{reflect}

\hypertarget{assessment-14}{%
\subsection*{Assessment}\label{assessment-14}}
\addcontentsline{toc}{subsection}{Assessment}

See the Assessment section in Moodle for assignment details and due dates.

\hypertarget{resources-7}{%
\subsection*{Resources}\label{resources-7}}
\addcontentsline{toc}{subsection}{Resources}

Here are the resources you will need to complete this unit.

\begin{itemize}
\tightlist
\item
  Chapter 8 of course text: \emph{Digital Filmmaking: A Beginner's Guide to Mastering the Craft,} by Ned Vankevich (e-text)
\item
  Other resources will be available online and in the course text.
\end{itemize}

\hypertarget{types-of-film-and-genres}{%
\section{Types of Film and Genres}\label{types-of-film-and-genres}}

{Types of Film }

There are three main types of short films that we can do for the project ahead:

\begin{itemize}
\tightlist
\item
  Classical Narrative
\item
  Documentary
\item
  Experimental, \emph{Avant-Garde}, and Surreal Cinema
\end{itemize}

Most of you will probably choose the classical narrative, but you are free to do the other two with the caveat that experimental/\emph{avant-garde} films are far more challenging than they might appear.

As a transferable skill, like in many industries and businesses, smart filmmakers look for a market before they commit too much time and money to a film project. Others of course rely on their intuition and passion. Either way, it is wise to see what those who distribute or stream and screen short films, such as festivals and competitions, look for in making their choices. Their criteria can you vet the quality of your story and its execution. See Film Shortage's \href{https://filmshortage.com/what-we-like/}{What Type Of Short Films Do We Prefer?}

{Film Genres}

The main films genres include:

\begin{itemize}
\tightlist
\item
  Drama\\
\item
  Comedy\\
\item
  Romance\\
\item
  Action\\
\item
  Thriller\\
\item
  Horror\\
\item
  Gangster\\
\item
  Crime\\
\item
  Adventure\\
\item
  Westerns\\
\item
  Sci-Fi\\
\item
  Fantasy\\
\item
  Historical\\
\item
  Epic\\
\item
  War\\
\item
  Bio Pics (biographies)
\end{itemize}

Knowing your film genre is important because it embodies narrative patterns and expectations audiences have, and if these are not fulfilled you will lessen your impact on those watching your film. Put another way, comedies need to be funny, horror films need to induce fear, romantic films need to inspire us to love, etc. Commercially, if these basic requirements are not met it can result in financial disaster and loss of reputation. At your level it will mean your film is not as effective as it could be.

Knowing your genre will also help you develop your story and script in this and the following unit.

\hypertarget{activity-planning-your-final-film-project}{%
\subsection*{Activity: Planning your Final Film Project}\label{activity-planning-your-final-film-project}}
\addcontentsline{toc}{subsection}{Activity: Planning your Final Film Project}

\begin{reflect}
Log in your journal what type of film you want to work on for your final film project and why.

Next, view this list of loglines (short one line descriptions of a film story) in various genres: \href{https://www.slideshare.net/BigBadBoobyLady/popular-short-film-genres}{Popular Short film genres} Which genre appeals to you most? Why? Log this reflection and start to brainstorm genre-based story concepts.
\end{reflect}

\hypertarget{logistics-for-the-final-film-project}{%
\section{Logistics for the Final Film Project}\label{logistics-for-the-final-film-project}}

This last film is a major project and will account for a lot of your grade given that it will reveal what you learned during this course. In light of this, you will need to create a strong story and the ``blueprint'' or script for filming it.

To save time so that you can put more of your energy into making the film, you will not have to develop a full screenplay (hence the scriptment section below) for the final project, but you are strongly encouraged to do so.

However, this does not mean you do not have to pitch a core story concept (a basic story, its genre, and hook or unique angle) to the class. Everyone will be required to do this so that you can gain more experience with your creativity and your pitching skills.

Once you have created a story concept that inspires you, pitch it to the class to see if other classmates want to join you in your project and to get feedback on it. Once you have your team in place you will then follow the process for creating a shooting script for pre-production and production.

\hypertarget{activity-story-research}{%
\subsection*{Activity: Story Research}\label{activity-story-research}}
\addcontentsline{toc}{subsection}{Activity: Story Research}

\begin{reflect}
In the next activity, you'll be asked to brainstorm and develop a core story idea to pitch to the class.

To help you do this, follow the steps in Story Research Section of Chapter Eight and ask and answer the questions in the What to Look for Section. This should prime your creative pump for the next topic. Also view this helpful resource:

\href{https://www.youtube.com/watch?v=wMqIQcTMlA0}{Watch: How to Write a Short Film}
\end{reflect}

\hypertarget{activity-brainstorm-and-prepare-your-pitch}{%
\subsection*{Activity: Brainstorm and Prepare Your Pitch}\label{activity-brainstorm-and-prepare-your-pitch}}
\addcontentsline{toc}{subsection}{Activity: Brainstorm and Prepare Your Pitch}

\begin{reflect}
Brainstorm either alone or with a partner and come up with a story concept to pitch to the class. The goal of this activity is to do find a compelling story and to pitch it in a way that will attract the teammates you will need to make your longer film.

Check the \href{assets/unit8/Guidelines_for_fairytale_template_for_a_short_film.pdf}{Guidelines for Pitching a Film Project}

To prepare for your pitch make sure you:

\begin{itemize}
\tightlist
\item
  Know your genre. The story above could be a comedy, a drama, thriller, or a horror film. Which genre you decide will determine the following.\\
\item
  Establish your main characters.\\
\item
  Establish the main conflict.\\
\item
  Escalate the conflict.\\
\item
  Create some of the obstacles that will have to be overcome.\\
\item
  Have a clear beginning, middle, and end.\\
\item
  Know what crew members you will need and what they will do, e.g., a screenwriter, producer, cinematographer, editor, etc.
\end{itemize}

Make sure the story is compressed, a la Aristotle's \emph{Poetics}, with a limited number of characters, locations, and action.

\textbf{Helpful Hints:}

\begin{itemize}
\tightlist
\item
  Sometimes a title for your film can lend lots of inspiration. Professional screenwriters and Hollywood spend a lot of time finding the right title that captures the spirit of the film and can aid in marketing. Think \emph{Rebel Without A Cause, Slum Dog Millionaire, Vertigo, Dumb and Dumber, The Haunting, Groundhog Day, When Harry Met Sally.}
\item
  Most of the short films that won or were nominated for the Academy Award for Best Live Action Short Film follow Aristotle's idea of great drama and comedy: a few characters and locations and action that takes place over a brief amount of time.
\end{itemize}
\end{reflect}

\hypertarget{creating-a-scriptment}{%
\section{Creating a Scriptment}\label{creating-a-scriptment}}

Our goal in this Unit is to find and create a strong story to film to shoot and edit so that it shines with excellence. The following steps will help you prepare a shooting script that will make pre-production, production, and post-production flow better.

\begin{itemize}
\tightlist
\item
  Create a step outline.\\
\item
  Create a scriptment.
\end{itemize}

Your scriptment should be written in a way that someone who reads it should be able to visualize and hear your film.

\hypertarget{activity-creating-a-scriptment}{%
\subsection*{Activity: Creating a Scriptment}\label{activity-creating-a-scriptment}}
\addcontentsline{toc}{subsection}{Activity: Creating a Scriptment}

\begin{reflect}
For an overview of scriptments, consult this resource: \href{https://screencraft.org/blog/outlines-treatments-and-scriptments-oh-my/}{Outlines, Treatments, and Scriptments, Oh My!} by Ken Miyamoto.

As you do other drafts (good writing involves a lot of re-writing) to improve your story and polish your script so that everything is clear, it is helpful to ask the following questions. (These questions are designed for a fiction narrative film, though some of them can help vet your documentary or experimental film.) You might not have answers to all of them but you should for most of them. Again the emphasis in on helping you to be clear about your project.

\begin{itemize}
\tightlist
\item
  Do I have the right genre?\\
\item
  Do I have a hook or is there something unique and fascinating about my story? It's special angle.\\
\item
  Is the story clear? A solid beginng, middle, and end?\\
\item
  Do I focus on only a few characters and limited locations?\\
\item
  Are my action and conflict unified and focused?\\
\item
  Does my main character have a goal?\\
\item
  Is there an antagonist who or which wants to thwart that goal? (Remember, weather or a dog can be an antagonist.)\\
\item
  Are their stakes or consequences for not obtaining the goal?\\
\item
  Are there interesting obstacles to obtaining the goal?\\
\item
  Do I have a set-up of the main character, problem, goal, and antagonist?\\
\item
  Do I have some unexpected and unpredictable twists and turns in the middle section of my story?\\
\item
  Is there a main conflict or battle that determines whether the main character gets his or her goal?\\
\item
  Do I have a satisfying ending? (Remember not all endings have to be happy or closed. Open-ended endings are not resolved and make the audience guess, debate, and discuss what happened or might happen after).
\end{itemize}

Once you have your team and have developed the scriptment for your story, you can begin pre-production to get everything in place to film your project: e.g., cast, locations, props, gear, shooting schedule, etc.
\end{reflect}

\hypertarget{summary-7}{%
\section*{Summary}\label{summary-7}}
\addcontentsline{toc}{section}{Summary}

In this unit, you learned about:

\begin{itemize}
\tightlist
\item
  Types of short films\\
\item
  Film Genres\\
\item
  The Logistics for the Final Project\\
\item
  Script and Story Development\\
\item
  Creating A Scriptment
\end{itemize}

\hypertarget{assessment-15}{%
\section*{Assessment}\label{assessment-15}}
\addcontentsline{toc}{section}{Assessment}

\begin{assessment}
{Course Journal}

After completing this unit, including the learning activities, you are asked to make sure you are doing journal entries and when appropriate to share your responses with your facilitator and classmates when you meet.

Also, log in your journal what you learned from the creating and filming of your In-Camera exercise. Log what you thought about the films of your peers and what you learned from the feedback from them and your instructor.

Be sure to make a note of what films stood out and who did them. You will consult this at the end of the course when the class determines: Best Overall Film. Most Imaginative Film, Best Story, Best Cinematography, Best Editing, etc.

\emph{See the Assessments section for more details on submitting your journal, as well as the grading criteria.}
\end{assessment}

\hypertarget{checking-your-learning-7}{%
\section*{Checking your Learning}\label{checking-your-learning-7}}
\addcontentsline{toc}{section}{Checking your Learning}

\begin{progress}
Before you move on to the next unit, you may want to check to make sure that you are able to:

\begin{itemize}
\tightlist
\item
  Describe the different types of short films\\
\item
  Define film genre and why it is important\\
\item
  Determine how to approach making a film script to shoot and edit.\\
\item
  Create a shooting script.
\end{itemize}
\end{progress}

\hypertarget{making-your-short-film}{%
\chapter{Making Your Short Film}\label{making-your-short-film}}

\hypertarget{overview-8}{%
\section*{Overview}\label{overview-8}}
\addcontentsline{toc}{section}{Overview}

In this unit we culminate this course as you make your final film project and bring together all you have learned so far. Congratulations. Your hard work is paying off.

Here you will get the other skills needed to make an excellent short film as you take your scriptment from the last unit and shoot and edit it.

These skills will focus on exploring how to use film coverage to create meaning, emotions, and moods.

Many of the items and techniques addressed here have been covered in the previous units. Go back to them if you need a refresher. But there are important new elements that you will learn. The important thing is not to get lost in the detail but to maintain the big picture perspective, namely, communicating what you want to say in a creative and excellent way.

Let's dig in.

\hypertarget{topics-8}{%
\subsection*{Topics}\label{topics-8}}
\addcontentsline{toc}{subsection}{Topics}

This unit is divided into the following topics:

\begin{itemize}
\tightlist
\item
  The Syntax of Film
\item
  The Importance and Types of Camera Coverage
\item
  Continuity
\item
  Storyboarding and Shot Lists
\item
  Casting
\item
  Sound Recording
\item
  The Final Steps
\end{itemize}

\hypertarget{learning-outcomes-8}{%
\subsection*{Learning Outcomes}\label{learning-outcomes-8}}
\addcontentsline{toc}{subsection}{Learning Outcomes}

When you have completed this unit, you should be able to:

\begin{itemize}
\tightlist
\item
  Define film syntax and its key terms to describe your filmmaking process.\\
\item
  Describe camera coverage and contrast the types of Coverage.\\
\item
  Demonstrate How To Use Continuity.\\
\item
  Create Storyboards and Shot Lists.\\
\item
  Demonstrate Casting Skills.\\
\item
  Apply L Cuts and J Cuts.
\end{itemize}

\hypertarget{activity-checklist-8}{%
\subsection*{Activity Checklist}\label{activity-checklist-8}}
\addcontentsline{toc}{subsection}{Activity Checklist}

\begin{reflect}
Here is a checklist of learning activities you will benefit from in completing this unit. You may find it useful for planning your work.

\begin{itemize}
\tightlist
\item
  Study the next section of Chapter 9 and watch the videos related to it. The selected videos for this unit are extremely important for you to understand how to make your final project. Failure to study them could jeopardize your final project.
\item
  Please log in your journal your key takeaways from the watched videos along with your responses to specific questions.
\item
  Storyboard the scenes from your scriptment you plan to film.
\item
  Review Unit 7 to refresh your knowledge about what is involved in the post-production process.
\item
  Study the guidelines and the Evaluation Criteria closely for the final project. They will help you to do well.
\end{itemize}
\end{reflect}

\hypertarget{assessment-16}{%
\subsection*{Assessment}\label{assessment-16}}
\addcontentsline{toc}{subsection}{Assessment}

See the Assessment section in Moodle for assignment details and due dates.

\hypertarget{resources-8}{%
\subsection*{Resources}\label{resources-8}}
\addcontentsline{toc}{subsection}{Resources}

Here are the resources you will need to complete this unit.

\begin{itemize}
\tightlist
\item
  Chapter 8 of course text: \emph{Digital Filmmaking: A Beginner's Guide to Mastering the Craft,} by Ned Vankevich (e-text)
\item
  Other resources will be available online and in the course text.
\end{itemize}

\hypertarget{the-syntax-of-film}{%
\section{The Syntax of Film}\label{the-syntax-of-film}}

The first principle of this unit will focus on film syntax, where you take the grammar of film---the types of shots, shot sizes, shot angles and height, shot framing, subject angles, etc. and order them in a way that best communicates what you want to say. This is where higher-level creativity takes place.

Filmmaking is like a language where you use the elements of grammar to make sentences. In this case, your shots to make scenes and the meaning the scenes will covey. The cumulation of sentences lead to paragraphs -- the film equivalent of sequences, and the cumulation of paragraphs leads to sections (acts in film) and the cumulation of sections leads to your essay (in this case, your final film). You get the metaphor.

\hypertarget{activity-read-watch-and-reflect}{%
\subsection*{Activity: Read, Watch and Reflect}\label{activity-read-watch-and-reflect}}
\addcontentsline{toc}{subsection}{Activity: Read, Watch and Reflect}

\begin{reflect}
Read and study the first section of Chapter 9: The Syntax of Film. Then, before moving on to the next topic in this unit, watch the following tutorials to re-enforce the big picture. Do not forget to log in your journal your best take-aways from them.

\href{https://www.youtube.com/watch?v=IK2IAEO-FUI}{Watch: How to Shoot a Scene! - Film Riot}

\begin{center}\rule{0.5\linewidth}{0.5pt}\end{center}

\href{https://www.youtube.com/watch?v=9AGaECt9j4g}{Watch: Film Blocking Tutorial --- Filmmaking Techniques for Directors: Ep3}
\end{reflect}

\hypertarget{the-importance-and-types-of-camera-coverage}{%
\section{The Importance and Types of Camera Coverage}\label{the-importance-and-types-of-camera-coverage}}

Exploring and understanding camera coverage is vital for those who want to create excellent films. Coverage refers to how a scene is captured. It involves how many shots are used and their type and kind to capture a scene in a film or video production. Coverage is thus the shot material an editor or post-production team will use to assemble the scenes and sequences of a movie. Having proper coverage is critical if a scene and a film is to make sense.

The are many questions that can guide you as you break down your film script to determine what coverage to use and how to shoot a scene. (Remember a scene can be one shot such as driving to the store, or a series of shots such as someone confronting a person to get information from him or her.) Take the time to study the \textbf{Questions for Discovering Coverage} section.

Coverage can be done well or poorly. Our goal is the former.

There are two main types of general coverage: \textbf{Master Shots and Mini-Masters}. Knowing why these are important and when to use them will help guide you to the more detailed coverage and shots you will employ as as you develop the action of a scene.

Filmmaking is not a ``paint by numbers'' art. If you follow formulas, most likely your film will be predicable and lack freshness. This is why understanding shot progression is important. Choosing the types of shots you will use and their order will form a large part of your visual creativity in your film.

\hypertarget{activity-read-and-watch-1}{%
\subsection*{Activity: Read and Watch}\label{activity-read-and-watch-1}}
\addcontentsline{toc}{subsection}{Activity: Read and Watch}

\begin{reflect}
Study the next section of Chapter 9 ``The Importance and Types of Camera Coverage''. Next, watch the video \href{https://www.youtube.com/watch?v=cz3nBkIa9K0}{What is a Master Shot}

\begin{center}\rule{0.5\linewidth}{0.5pt}\end{center}

{Questions to Consider}

\begin{itemize}
\tightlist
\item
  Explain the value of why this type of shot can be beneficial as the first shot you do when you start your coverage of a scene? What is the downside of relying too much on a master shot?
\end{itemize}
\end{reflect}

\hypertarget{continuity}{%
\section{Continuity}\label{continuity}}

Continuity, or the non-distracting and motivated and smooth flow of action, is another skill set essential to the filmmaker's took kit. Some techniques you have already studied, but others are new and important.

The following are the major techniques of continuity that you must focus on:

\begin{itemize}
\tightlist
\item
  Matching Action\\
\item
  Directional Continuity\\
\item
  Eyeline Continuity\\
\item
  180-Degree Rule\\
\item
  30-Degree Rule
\end{itemize}

Each of these individually and as a group will help immensely in creating well-constructed, well-motivated, and audience-engaging films.

\hypertarget{activity-read-and-watch-2}{%
\subsection*{Activity: Read and Watch}\label{activity-read-and-watch-2}}
\addcontentsline{toc}{subsection}{Activity: Read and Watch}

\begin{reflect}
Read the Continuity section of Chapter 9.

There are many resources that can help you to understand the various parts of continuity. It might seem like a lot but it will greatly benefit you if you view the following resources.

\begin{itemize}
\tightlist
\item
  \href{https://www.youtube-nocookie.com/embed/eou7A-e2e4I}{Match on Action technique}
\item
  \href{https://www.youtube-nocookie.com/embed/El28XrjtcMI}{Match Cuts in Film Editing}
\item
  \href{https://www.youtube-nocookie.com/embed/9XOn5uxdSJc}{Screen Direction rule}
\item
  \href{https://www.youtube-nocookie.com/embed/RogoUz_pk4Y}{Screen direction}
\item
  \href{https://www.youtube-nocookie.com/embed/HinUychY3sE}{Breaking Down the 180-Degree Rule}
\item
  \href{https://www.youtube-nocookie.com/embed/1K8EUc98VoQ}{The 30 Degree Rule in Filmmaking \textbar{} Kaicreative \textbar{} Tips}
\end{itemize}

Do you feel you have a stronger grasp of continuity and of each the techniques addressed? If not, review the videos of those parts you do not understand.
\end{reflect}

\hypertarget{storyboarding}{%
\section{Storyboarding}\label{storyboarding}}

By now you might feel overwhelmed with camera coverage and all that it takes to shoot your story and script well. Storyboarding and storyboards can help simplify the process and make it easier to visualize what you want to do and how to do it.

It will help your scene coverage if you storyboard your shots. They do not have to be elaborate. They can be simple stick figures such as this one that shows a long shot of someone pushing or struggling against something (the context of your story will determine this).

\begin{figure}
\centering
\includegraphics{assets/unit5/Picture1.png}
\caption{\emph{graphic stick figure pushing}}
\end{figure}

The important thing is not the quality but that your storyboards make sense to you and your crew so you will know how to approach a shot and why.

At the same time, do not become a slave to your storyboard. If you discover something better or get inspired on the film set, try the new thing. For safety, shoot it the storyboard way and then the new way and decide during post-production which one will work best.

\hypertarget{shot-lists}{%
\section{Shot Lists}\label{shot-lists}}

Once you have your scriptment describing all of the action and dialogue you will film and have storyboarded your film, you can now create a shot list which will help ensure you get all the shots needed for your coverage.

It is better if your whole production team gives feedback on the shot list. Ultimately, it is the director who will make the final choice but hearing what other team members feel can help cut down on mistakes.

Study the Shot List Example in Chapter Nine to see what you need to include.

\hypertarget{activity-read-and-watch-3}{%
\subsection*{Activity: Read and Watch}\label{activity-read-and-watch-3}}
\addcontentsline{toc}{subsection}{Activity: Read and Watch}

\begin{reflect}
\begin{itemize}
\tightlist
\item
  Read the section on Storyboarding and Shot Lists in Chapter 9.
\item
  Next, read the article \href{https://www.wyzowl.com/what-is-a-storyboard/}{``What Is a Storyboard and How Do You Make One for Your Video?} to better understand storyboards
\item
  Storyboard the scenes from your scriptment you plan to film. You can do this alone but it will be better if at least the director and cinematographer for your project work on this together. Ideally everyone on the team should work on this so that you will bond and have a unified vision of how to stage and shoot your shots.
\item
  Watch the video \href{https://www.youtube.com/watch?v=-rzJP_5L_yQ}{``Master the Art of Shot Lists: Boost Your Filmmaking Efficiency!''}
\end{itemize}

\begin{itemize}
\tightlist
\item
  Make sure to include your shot lists in your final project and course journal.
\end{itemize}

\textbf{Helpful Hint:} Remember your locations will determine a lot of how you can and cannot shoot a scene. A drone will not work in a small room. Therefore, be sure to scout your locations and take pictures of them as you desing your storyboards and shot lists. Be sure to also confirm that the locations will be available for your shooting schedule.
\end{reflect}

\hypertarget{casting}{%
\section{Casting}\label{casting}}

Once you have your locations, storyboards, shot list, and schedule you are ready to cast your show. Casting is often done earlier in larger professional productions because they need to lock in busy actors. For this scale of micro or no budget filmmaking having your shooting schedule and your storyboards and shot lists can show seasoned actors that you know what you are doing and can help you land them.

This pre-production process emphasized here will also help you to know if some special skills are needed for your actors. E.g., someone who can dance, sing, play guitar, is good at soccer, etc.

Having good actors is critical to having a good film. Follow the tips and guidelines in the chapter and work hard to find the best actors you can. It will pay off a lot.

\hypertarget{read-and-watch}{%
\subsection*{Read and Watch}\label{read-and-watch}}
\addcontentsline{toc}{subsection}{Read and Watch}

\begin{reflect}
Read the Casting section in Chapter 9.

Before auditioning and casting your film, watch the video \href{https://www.youtube.com/watch?v=YpCkRPqsiJ4}{``How to Cast an ACTOR for a No Budget Film''} for a simple overview of a casting process

For practical tips for auditioning actors watch the video \href{https://www.youtube.com/watch?v=x0G6n346m90}{``Auditioning Actors''}

For detailed information regarding how to work with actors and casting for low budget films listen to this podcast:\href{https://www.youtube.com/watch?v=-DwKilT0T34}{How to Cast a No Budget Indie Film with Casting Director Veronika Lee}
\end{reflect}

\hypertarget{sound-recording}{%
\section{Sound Recording}\label{sound-recording}}

Audio can make or break a film, especially in no-budget and low-budget filmmaking. Proper sound recording and sound editing is a course unto itself and if you want to go into filmmaking professionally you should take at least one course in this.

\hypertarget{activity-read-and-research}{%
\subsection*{Activity: Read and Research}\label{activity-read-and-research}}
\addcontentsline{toc}{subsection}{Activity: Read and Research}

\begin{reflect}
How to capture great sound for a film without the use of professionals and professional gear is a challenge. View the three resources for this in Chapter Nine and find tutorials and instruction essays on the web that are appropriate to your cell phone and DSLR.
\end{reflect}

\hypertarget{the-final-steps}{%
\section{The Final Steps}\label{the-final-steps}}

Bravo. Once you have reached this stage of this Unit and the previous one you should be well-equipped to produce and film your movie.

You will then have to edit it. This section adds some new techniques that can help elevate your film project to a higher level, namely Split Edits. If you can, add this to the editing of your film, as well as L cuts and J cuts. They will make your film flow better and be more engaging and enjoyable. It might also help win you an award in the course.

\hypertarget{activity-review}{%
\subsection*{Activity: Review}\label{activity-review}}
\addcontentsline{toc}{subsection}{Activity: Review}

\begin{reflect}
Review Unit 7 to refresh yourself about what is involved in the post-production process. If possible, add Split Edits to elevate the quality and sophistication of your final film project.
\end{reflect}

\hypertarget{activity-guidelines-for-doing-and-evaluating-the-final-film-project}{%
\subsection*{Activity: Guidelines for Doing and Evaluating The Final Film Project}\label{activity-guidelines-for-doing-and-evaluating-the-final-film-project}}
\addcontentsline{toc}{subsection}{Activity: Guidelines for Doing and Evaluating The Final Film Project}

\begin{reflect}
Continuity or the non-distracting and motivated and smooth flow of action is another skill set essential to the filmmaker's toolkit. Study the following resources and practice applying them in your edited scenes especially those with dialogue or where one scene transitions into another:

\begin{itemize}
\tightlist
\item
  \href{https://www.youtube-nocookie.com/embed/eyH-a964kAs}{SFX Secrets: The J Cut \& The L Cut}
\end{itemize}

\begin{center}\rule{0.5\linewidth}{0.5pt}\end{center}

\begin{itemize}
\tightlist
\item
  \href{https://www.youtube-nocookie.com/embed/fT5rRPMnka0}{Video Editing Tips: J Cut vs L Cut}
\end{itemize}

In addition, study the \emph{Steps to Making an Excellent Short Film} section in the chapter. It is both a good recap and a good checklist to help ensure you have the elements need to make an excellent final film.

Please study the guidelines and the Evaluation Criteria closely for the final project. They will help you to do well. This criteria closely tracks that which was used to evaluate your montages so the process should be more comfortable for you now.

\begin{itemize}
\tightlist
\item
  Did the film make sense and work overall?\\
\item
  Did the film meet genre expectations?\\
\item
  Was the story interesting? Surprising, engaging, exciting?\\
\item
  Did the actors perform well?\\
\item
  Did the camera coverage work? If so, why? If not, why?\\
\item
  Did the editing work? If so why? If not, why?\\
\item
  What stood out in a good way?\\
\item
  What needed improvement?\\
\item
  What did you like best about the film?\\
\item
  What did you like least?\\
\item
  Other comments.
\end{itemize}
\end{reflect}

\hypertarget{summary-8}{%
\section*{Summary}\label{summary-8}}
\addcontentsline{toc}{section}{Summary}

In this unit, you learned about:

\begin{itemize}
\tightlist
\item
  The Syntax of Film\\
\item
  The Importance and Types of Coverage\\
\item
  Continuity\\
\item
  Storyboarding and Shot Lists\\
\item
  Casting\\
\item
  Sound Recording
\end{itemize}

\hypertarget{assessment-17}{%
\section*{Assessment}\label{assessment-17}}
\addcontentsline{toc}{section}{Assessment}

\begin{assessment}
{Course Journal}

After completing this unit, including the learning activities, you are asked to make sure you are doing journal entries and when appropriate to share your responses with your facilitator and classmates when you meet.

Note that entries are expected after every unit. Your journal reflections are submitted at the end of the course as part of the Journal One: Personal Journal and self assessment.

\emph{See the Assessments section for more details on submitting your journal, as well as the grading criteria.}

{Short Film Project }

This final short film project will account for 40\% of your grade. See the Assessments section for more details on your Final short Film Project and how to submit it, as well as the grading criteria.
\end{assessment}

\hypertarget{checking-your-learning-8}{%
\section*{Checking your Learning}\label{checking-your-learning-8}}
\addcontentsline{toc}{section}{Checking your Learning}

\begin{progress}
Before you move on to the next unit, you may want to check to make sure that you are able to:

\begin{itemize}
\tightlist
\item
  Define film syntax and its key terms to describe your filmmaking process.\\
\item
  Describe Camera Coverage and Contrast The Types of Coverage.\\
\item
  Demonstrate How To Use Continuity.\\
\item
  Create Storyboards and Shot Lists.\\
\item
  Demonstrate Casting Skills.\\
\item
  Apply L Cuts and J Cuts.
\end{itemize}
\end{progress}

  \bibliography{book.bib}

\end{document}
